% =====================================================================
%  The Idiot Brain - shared Beamer preamble for all chapter decks
% ---------------------------------------------------------------------
%  Usage (in slides/chNN/chNN.tex):
%      \documentclass[aspectratio=169,12pt]{beamer}
%      % =====================================================================
%  The Idiot Brain webinars - shared Beamer preamble (the banner look)
% ---------------------------------------------------------------------
%  Usage (in slides/webinarNN/webinarNN.tex):
%      \documentclass[aspectratio=169,11pt]{beamer}
%      % =====================================================================
%  The Idiot Brain webinars - shared Beamer preamble (the banner look)
% ---------------------------------------------------------------------
%  Usage (in slides/webinarNN/webinarNN.tex):
%      \documentclass[aspectratio=169,11pt]{beamer}
%      % =====================================================================
%  The Idiot Brain webinars - shared Beamer preamble (the banner look)
% ---------------------------------------------------------------------
%  Usage (in slides/webinarNN/webinarNN.tex):
%      \documentclass[aspectratio=169,11pt]{beamer}
%      \input{../common/preamble.tex}
%  Build from the deck folder:  latexmk -lualatex -outdir=build webinarNN.tex
%  Art (emoji, mascot moods, title banner) comes from common/build_art.py.
%
%  All styling lives here, so decks only hold content. Palette and fonts
%  are the banner's (banner/build_banner.py): plum text on cream, pastel
%  accents, Mardoto for text, Glazz for big titles, Twemoji icons and
%  the pink brain mascot (DECISIONS.md #4-#6).
% =====================================================================

% ---------------------------------------------------------------------
% Fonts (LuaLaTeX + fontspec), shared with the banner
% ---------------------------------------------------------------------
\usepackage{fontspec}
\tracinglostchars=3 % a glyph missing from a font stops the build instead of printing nothing
\defaultfontfeatures{Ligatures=TeX}
\def\idiotfontdir{../../banner/assets/fonts/}
\setsansfont{Mardoto}[Path=\idiotfontdir, Extension=.ttf, UprightFont=*-Regular, BoldFont=*-Bold]
\setmainfont{Mardoto}[Path=\idiotfontdir, Extension=.ttf, UprightFont=*-Regular, BoldFont=*-Bold]
\newfontfamily\mediumfont{Mardoto-Medium}[Path=\idiotfontdir, Extension=.ttf]
\newfontfamily\blackfont{Mardoto-Black}[Path=\idiotfontdir, Extension=.ttf]
\newfontfamily\titlefont{GlazzArm-Regular}[Path=\idiotfontdir, Extension=.ttf] % Armenian letters only
\newfontfamily\handfont{Ink Free} % Windows font with Armenian; speech bubbles only
% Mardoto has no italic: emphasis is medium weight in the accent colour instead.
\renewcommand{\emph}[1]{{\mediumfont\color{plum}#1}}
% No hyphenation: split words are hard to read on a shared screen.
\hyphenpenalty=10000
\exhyphenpenalty=10000

% ---------------------------------------------------------------------
% Packages
% ---------------------------------------------------------------------
\usepackage{graphicx}
\usepackage{booktabs}
\usepackage[most]{tcolorbox}
\usepackage{tikz}
\usetikzlibrary{arrows.meta, positioning, calc, shapes.geometric, shapes.callouts, shapes.misc,
  decorations.pathmorphing, decorations.pathreplacing, backgrounds, fit, patterns.meta}
\usepackage{pgfplots}
\pgfplotsset{compat=1.18}
\usepgfplotslibrary{fillbetween}
\graphicspath{{../common/art/}{images/}}

% ---------------------------------------------------------------------
% Colours: the banner palette. Charts with 3+ series use the Armenian
% flag colours (amRed / amBlue / amOrange).
% ---------------------------------------------------------------------
\definecolor{plum}{HTML}{4A2340}
\definecolor{brainpink}{HTML}{F7A8C4}
\definecolor{fold}{HTML}{E57FA6}
\definecolor{cream}{HTML}{FFF6EC}
\definecolor{cardwhite}{HTML}{FFFDF9}
\definecolor{peach}{HTML}{FFD9C2}
\definecolor{lavender}{HTML}{D9C8F5}
\definecolor{mint}{HTML}{B5E8D5}
\definecolor{butter}{HTML}{FFE8A3}
\definecolor{blush}{HTML}{FF8FA8}
\definecolor{amRed}{HTML}{D90012}
\definecolor{amBlue}{HTML}{0033A0}
\definecolor{amOrange}{HTML}{F2A800}
% darker partners for text on cream
\colorlet{mintdark}{mint!45!black!80!green}
\definecolor{mintdark}{HTML}{2E7D62}
\definecolor{lavenderdark}{HTML}{5B3FA0}
\definecolor{peachdark}{HTML}{B5532A}
\definecolor{pinkdark}{HTML}{C2185B}
% names used by the figure files
\colorlet{ink}{plum}
\colorlet{paper}{cream}
\colorlet{mist}{lavender!40!cream}
\definecolor{night}{HTML}{2B2150}
\colorlet{softgray}{plum!60!cream}
% accent of the current part; \partaccent{<colour>} switches it
\colorlet{accent}{butter}
\newcommand{\partaccent}[1]{\colorlet{accent}{#1}}

% ---------------------------------------------------------------------
% Art: \emoji, \mascot (pages of common/art/*.pdf, see build_art.py)
% ---------------------------------------------------------------------
\makeatletter
\newcommand{\artpage}[3]{\expandafter\def\csname idiot@art@#1@#2\endcsname{#3}}
\input{../common/art/art_index.tex}
\newcommand{\idiot@art}[3]{% kind, name, size options
  \ifcsname idiot@art@#1@#2\endcsname
    \includegraphics[#3, page=\csname idiot@art@#1@#2\endcsname]{#1.pdf}%
  \else
    \PackageError{idiot}{Unknown #1 '#2'; add it to build_art.py and rebuild the art}{}%
  \fi}
% \emoji[<height>]{name}: inline icon, sits on the text baseline like a letter
\newcommand{\emoji}[2][1.15em]{\raisebox{-0.2em}{\idiot@art{emoji}{#2}{height=#1}}}
% \bigemoji[<height>]{name}: no baseline shift, for TikZ nodes and big pictures
\newcommand{\bigemoji}[2][1.5cm]{\idiot@art{emoji}{#2}{height=#1}}
% \mascot[<width>]{mood}: panic, happy, scared, sleepy, sick, confused, smug, hungry, thinking
\newcommand{\mascot}[2][2.6cm]{\idiot@art{mascot}{#2}{width=#1}}
\makeatother
% the older figure files place the mascot with \pic{noggin=<mood>}
\tikzset{pics/noggin/.style={code={\node[inner sep=0, transform shape] {\mascot[2.7cm]{#1}};}}}

% ---------------------------------------------------------------------
% Beamer look
% ---------------------------------------------------------------------
\setbeamertemplate{navigation symbols}{}
\setbeamercolor{normal text}{fg=plum, bg=cream}
\setbeamercolor{structure}{fg=plum}
\setbeamercolor{alerted text}{fg=pinkdark}
\setbeamerfont{alerted text}{series=\bfseries}
\setbeamercolor{frametitle}{fg=plum}
\setbeamersize{text margin left=9mm, text margin right=9mm}
\setbeamertemplate{frame begin}{\raggedright} % justified Armenian leaves big gaps

% cream paper with the banner's faint dot grid
\setbeamertemplate{background}{%
  \begin{tikzpicture}
    \fill[cream] (0,0) rectangle (\paperwidth,\paperheight);
    \fill[pattern={Dots[distance=4.4mm, radius=0.28mm]}, pattern color=plum!9!cream]
      (0,0) rectangle (\paperwidth,\paperheight);
  \end{tikzpicture}}

% topic pill, top right of every frame: set by \topic (emoji + short name)
\def\idiottopicemoji{}
\def\idiottopicname{}
\newcommand{\idiotpill}[3][cardwhite]{% fill, emoji, text
  \tikz[baseline=(p.base)]{\node[fill=#1, draw=plum, line width=0.7pt, rounded corners=2.6mm,
      inner xsep=2.2mm, inner ysep=1.1mm, font=\mediumfont\footnotesize, text=plum] (p)
      {\ifx\relax#2\relax\else\emoji[1.05em]{#2}\hspace{1.2mm}\fi#3};
    \begin{scope}[on background layer]
      \fill[plum, rounded corners=2.6mm] ([xshift=0.7mm, yshift=-0.7mm]p.south west) rectangle ([xshift=0.7mm, yshift=-0.7mm]p.north east);
    \end{scope}}}
% \chip{emoji}{text}: inline sticker label, same as the banner's topic chips
\newcommand{\chip}[3][cardwhite]{\idiotpill[#1]{#2}{#3}}

\setbeamertemplate{frametitle}{%
  \vskip 3.5mm
  \begin{minipage}[t]{0.66\paperwidth}
    \raggedright{\blackfont\fontsize{17.5}{20.5}\selectfont\insertframetitle}\par
    \vskip 1.3mm
    \tikz\fill[accent, rounded corners=0.6mm] (0,0) rectangle (1.5cm,1.1mm);
  \end{minipage}%
  \hfill
  \ifx\idiottopicname\empty\else
    \raisebox{-1mm}{\idiotpill{\idiottopicemoji}{\idiottopicname}}%
  \fi
  \vskip 1mm}

\setbeamertemplate{footline}{%
  \hbox to \paperwidth{\hskip 9mm
    {\mediumfont\fontsize{7}{8}\selectfont\color{plum!55!cream}Ապուշ ուղեղը · \insertshorttitle}\hfill
    {\blackfont\fontsize{8}{9}\selectfont\color{plum!70!cream}\insertframenumber}\hskip 9mm}
  \vskip 3mm}

\setbeamertemplate{itemize item}{\tikz[baseline=-0.65ex]\draw[plum, fill=accent, line width=0.6pt] (0,0) circle (0.95mm);}
\setbeamertemplate{itemize subitem}{\tikz[baseline=-0.65ex]\fill[plum!70] (0,0) circle (0.6mm);}
\setbeamertemplate{enumerate item}{\tikz[baseline=(n.base)]\node[fill=accent, draw=plum, line width=0.6pt,
  circle, inner sep=0.5mm, minimum size=4.2mm, font=\blackfont\scriptsize] (n) {\insertenumlabel};}
\setlength{\leftmargini}{5mm}
\setbeamertemplate{itemize/enumerate body begin}{\setlength{\itemsep}{1.6mm}}
\setbeamertemplate{itemize/enumerate subbody begin}{\setlength{\itemsep}{0.6mm}\small}

% ---------------------------------------------------------------------
% Structure: parts (full-colour openers) and topics (book sections)
% ---------------------------------------------------------------------
% \topic{<emoji>}{<short name>}: starts a book section; the pill shows on every following frame
\newcommand{\topic}[2]{\subsection{#2}\gdef\idiottopicemoji{#1}\gdef\idiottopicname{#2}}
\newcommand{\notopic}{\gdef\idiottopicemoji{}\gdef\idiottopicname{}}

% \partopener{<label>}{<title>}{<subtitle>}{<mood>}{<colour>}
\newcommand{\partopener}[5]{%
  \section{#2}\notopic\partaccent{#5}%
  \begin{frame}[plain]
    \begin{tikzpicture}[remember picture, overlay]
      \fill[#5] (current page.south west) rectangle (current page.north east);
      \fill[pattern={Dots[distance=4.4mm, radius=0.28mm]}, pattern color=plum!12!#5]
        (current page.south west) rectangle (current page.north east);
      \fill[cardwhite, opacity=0.45] ([xshift=-3.3cm, yshift=-0.2cm]current page.east) circle (3.1cm);
      \node[anchor=north west, inner sep=0] at ([xshift=10mm, yshift=-12mm]current page.north west)
        {\idiotpill{}{#1}};
      \node[anchor=west, text=plum, align=left, text width=9.2cm, inner sep=0]
        (t) at ([xshift=10mm, yshift=5mm]current page.west)
        {\titlefont\fontsize{38}{40}\selectfont #2};
      \node[anchor=north west, text=plum, text width=9.2cm, inner sep=0, align=left]
        at ([yshift=-5mm]t.south west) {\mediumfont\large #3};
      \node at ([xshift=-3.3cm, yshift=-0.2cm]current page.east) {\mascot[4.6cm]{#4}};
    \end{tikzpicture}
  \end{frame}}

% \topicopener{<emoji>}{<short name>}{<book's section title>}{<what it is about>}{<mood>}
\newcommand{\topicopener}[5]{%
  \topic{#1}{#2}%
  \begin{frame}[plain]
    \begin{tikzpicture}[remember picture, overlay]
      \fill[accent!55!cream] (current page.south west) rectangle (current page.north east);
      \fill[pattern={Dots[distance=4.4mm, radius=0.28mm]}, pattern color=plum!10!accent!55!cream]
        (current page.south west) rectangle (current page.north east);
      \node[anchor=north west, inner sep=0] at ([xshift=10mm, yshift=-12mm]current page.north west)
        {\idiotpill{#1}{#2}};
      \node[anchor=west, text=plum, text width=9.6cm, align=left, inner sep=0]
        (t) at ([xshift=10mm, yshift=6mm]current page.west)
        {\blackfont\fontsize{24}{28}\selectfont «#3»};
      \node[anchor=north west, text=plum, text width=9.6cm, align=left, inner sep=0]
        at ([yshift=-4mm]t.south west) {\mediumfont\large #4};
      \node at ([xshift=-3.0cm, yshift=-0.4cm]current page.east) {\mascot[4.2cm]{#5}};
      \node[rotate=-12] at ([xshift=-1.4cm, yshift=2.3cm]current page.east) {\bigemoji[1.5cm]{#1}};
    \end{tikzpicture}
  \end{frame}}

% ---------------------------------------------------------------------
% Boxes (sticker style: plum outline + hard plum shadow, like the banner)
% ---------------------------------------------------------------------
\tcbset{
  sticker/.style={enhanced, arc=3mm, boxrule=0.9pt, colframe=plum, colback=cardwhite, halign=left,
    left=3mm, right=3mm, top=2mm, bottom=2mm, fontupper=\normalsize, coltext=plum,
    shadow={1.1mm}{-1.1mm}{0mm}{fill=plum}},
  tagged/.style={sticker, colback=#1!45!cream, colbacktitle=#1, coltitle=plum,
    fonttitle=\blackfont\footnotesize, top=3mm,
    attach boxed title to top left={xshift=4mm, yshift*=-2.4mm},
    boxed title style={arc=2.2mm, boxrule=0.9pt, colframe=plum, left=1.5mm, right=1.5mm, top=0.6mm, bottom=0.6mm}},
}
\newtcolorbox{card}[1][]{sticker, #1}
\newtcolorbox{tryit}[1][]{tagged=mint, title={\emoji[1em]{testtube}\ Փորձիր հիմա}, #1}
\newtcolorbox{poll}[1][]{tagged=butter, title={\emoji[1em]{hand}\ Հարցում}, #1}
\newtcolorbox{chatbox}[1][]{tagged=lavender, title={\emoji[1em]{speech}\ Գրիր չաթում}, #1}
\newtcolorbox{scienceupdate}[1][]{tagged=peach, fontupper=\small, title={\emoji[1em]{microscope}\ Գիտությունն այսօր}, #1}
\newtcolorbox{bookquote}[1][]{sticker, colback=cardwhite, fontupper=\mediumfont, left=4mm,
  borderline west={1.4mm}{0pt}{accent},
  before upper={{\blackfont\footnotesize\emoji[1em]{book}\ Բըրնեթը գրում է}\par\smallskip}, #1}

% Brain.exe-style popup window: \begin{popup}{<window title>} ... \end{popup}
\newcommand{\winbuttons}{%
  \tikz[baseline=-0.6ex, line width=0.6pt]{
    \foreach \i/\f in {0/cardwhite, 1/cardwhite, 2/blush}{
      \filldraw[fill=\f, draw=plum, rounded corners=0.6mm] (\i*4.3mm,-1.6mm) rectangle ++(3.3mm,3.3mm);}
    \draw[plum] (0.8mm,-0.8mm) -- (2.5mm,-0.8mm);
    \draw[plum] (5.1mm,-0.8mm) rectangle (6.8mm,0.9mm);
    \draw[plum] (9.4mm,-0.9mm) -- (11.1mm,0.8mm) (9.4mm,0.8mm) -- (11.1mm,-0.9mm);}}
\newtcolorbox{popup}[2][]{sticker, colback=cardwhite, colbacktitle=brainpink, coltitle=plum,
  fonttitle=\blackfont\small, title={\emoji[1.1em]{warning}\hspace{1.5mm}#2\hfill\winbuttons},
  titlerule=0.9pt, toptitle=0.8mm, bottomtitle=0.8mm, #1}

% Myth or fact: \mythorfact{<statement>}{<verdict>}{<colour>}{<explanation>}; verdict shows on slide 2
\newcommand{\mythorfact}[4]{%
  \begin{center}
  \begin{tikzpicture}
    \node[draw=plum, line width=0.9pt, fill=cardwhite, rounded corners=4mm, inner xsep=8mm, inner ysep=6mm,
          text width=10cm, align=center, font=\blackfont\Large, text=plum] (mf) {«#1»};
    \begin{scope}[on background layer]
      \fill[plum, rounded corners=4mm] ([xshift=1.2mm, yshift=-1.2mm]mf.south west) rectangle ([xshift=1.2mm, yshift=-1.2mm]mf.north east);
    \end{scope}
    \node[anchor=west] at ([xshift=6mm]mf.north west) {\idiotpill[butter]{question}{Առասպել, թե՞ փաստ}};
    \onslide<2->{\node[draw=#3, text=#3, line width=2pt, rounded corners=2mm, rotate=-10, font=\blackfont\LARGE,
      inner sep=4pt, fill=cardwhite] at ([xshift=-8mm, yshift=1mm]mf.south east) {#2};}
  \end{tikzpicture}
  \end{center}
  \onslide<2->{#4}}

% Big question with the confused mascot: \bigquestion{<question>}{<small line>}
\newcommand{\bigquestion}[2]{%
  \vfill
  \begin{center}
  \begin{tikzpicture}
    \node[text width=9.6cm, align=center, font=\blackfont\fontsize{24}{29}\selectfont, text=plum] (q) {#1};
    \node[below=4mm of q, text width=9.6cm, align=center, font=\mediumfont\large, text=softgray] {#2};
    \node[anchor=east] at ([xshift=-2mm]q.west) {\mascot[2.6cm]{confused}};
  \end{tikzpicture}
  \end{center}
  \vfill}

% Small source line at the bottom of a frame: \source{Loftus & Palmer, 1974}
\newcommand{\source}[1]{%
  \begin{tikzpicture}[remember picture, overlay]
    \node[anchor=south east, font=\fontsize{6.5}{7.5}\selectfont, text=plum!60!cream, inner sep=0]
      at ([xshift=-17mm, yshift=3.4mm]current page.south east) {Աղբյուր՝ #1};
  \end{tikzpicture}}

% Speech bubbles (handwritten) and the overlay helper for TikZ:
% \node[visible on=<2->] ... keeps the layout but hides the node until slide 2.
\tikzset{
  say/.style={rectangle callout, draw=plum, line width=0.8pt, fill=cardwhite, rounded corners=3pt,
    inner sep=4pt, align=center, font=\handfont\normalsize, text=plum},
  think/.style={cloud callout, cloud puffs=12, draw=plum, line width=0.8pt, fill=cardwhite,
    inner sep=2pt, align=center, font=\handfont\normalsize, text=plum},
  % boxes for diagrams drawn directly in a frame: a style with #1 cannot be defined inside a frame body
  flow/.style={draw=plum, line width=0.8pt, fill=cardwhite, rounded corners=2mm, align=center,
    font=\small, text=plum, inner sep=4pt, text width=#1},
  flow/.default=3cm,
  chatmsg/.style={draw=plum, line width=0.8pt, rounded corners=3mm, inner sep=6pt, text width=#1,
    align=left, font=\small, text=plum},
  chatmsg/.default=7cm,
  invisible/.style={opacity=0, text opacity=0},
  visible on/.style={alt={#1{}{invisible}}},
  alt/.code args={<#1>#2#3}{\alt<#1>{\pgfkeysalso{#2}}{\pgfkeysalso{#3}}},
}

%  Build from the deck folder:  latexmk -lualatex -outdir=build webinarNN.tex
%  Art (emoji, mascot moods, title banner) comes from common/build_art.py.
%
%  All styling lives here, so decks only hold content. Palette and fonts
%  are the banner's (banner/build_banner.py): plum text on cream, pastel
%  accents, Mardoto for text, Glazz for big titles, Twemoji icons and
%  the pink brain mascot (DECISIONS.md #4-#6).
% =====================================================================

% ---------------------------------------------------------------------
% Fonts (LuaLaTeX + fontspec), shared with the banner
% ---------------------------------------------------------------------
\usepackage{fontspec}
\tracinglostchars=3 % a glyph missing from a font stops the build instead of printing nothing
\defaultfontfeatures{Ligatures=TeX}
\def\idiotfontdir{../../banner/assets/fonts/}
\setsansfont{Mardoto}[Path=\idiotfontdir, Extension=.ttf, UprightFont=*-Regular, BoldFont=*-Bold]
\setmainfont{Mardoto}[Path=\idiotfontdir, Extension=.ttf, UprightFont=*-Regular, BoldFont=*-Bold]
\newfontfamily\mediumfont{Mardoto-Medium}[Path=\idiotfontdir, Extension=.ttf]
\newfontfamily\blackfont{Mardoto-Black}[Path=\idiotfontdir, Extension=.ttf]
\newfontfamily\titlefont{GlazzArm-Regular}[Path=\idiotfontdir, Extension=.ttf] % Armenian letters only
\newfontfamily\handfont{Ink Free} % Windows font with Armenian; speech bubbles only
% Mardoto has no italic: emphasis is medium weight in the accent colour instead.
\renewcommand{\emph}[1]{{\mediumfont\color{plum}#1}}
% No hyphenation: split words are hard to read on a shared screen.
\hyphenpenalty=10000
\exhyphenpenalty=10000

% ---------------------------------------------------------------------
% Packages
% ---------------------------------------------------------------------
\usepackage{graphicx}
\usepackage{booktabs}
\usepackage[most]{tcolorbox}
\usepackage{tikz}
\usetikzlibrary{arrows.meta, positioning, calc, shapes.geometric, shapes.callouts, shapes.misc,
  decorations.pathmorphing, decorations.pathreplacing, backgrounds, fit, patterns.meta}
\usepackage{pgfplots}
\pgfplotsset{compat=1.18}
\usepgfplotslibrary{fillbetween}
\graphicspath{{../common/art/}{images/}}

% ---------------------------------------------------------------------
% Colours: the banner palette. Charts with 3+ series use the Armenian
% flag colours (amRed / amBlue / amOrange).
% ---------------------------------------------------------------------
\definecolor{plum}{HTML}{4A2340}
\definecolor{brainpink}{HTML}{F7A8C4}
\definecolor{fold}{HTML}{E57FA6}
\definecolor{cream}{HTML}{FFF6EC}
\definecolor{cardwhite}{HTML}{FFFDF9}
\definecolor{peach}{HTML}{FFD9C2}
\definecolor{lavender}{HTML}{D9C8F5}
\definecolor{mint}{HTML}{B5E8D5}
\definecolor{butter}{HTML}{FFE8A3}
\definecolor{blush}{HTML}{FF8FA8}
\definecolor{amRed}{HTML}{D90012}
\definecolor{amBlue}{HTML}{0033A0}
\definecolor{amOrange}{HTML}{F2A800}
% darker partners for text on cream
\colorlet{mintdark}{mint!45!black!80!green}
\definecolor{mintdark}{HTML}{2E7D62}
\definecolor{lavenderdark}{HTML}{5B3FA0}
\definecolor{peachdark}{HTML}{B5532A}
\definecolor{pinkdark}{HTML}{C2185B}
% names used by the figure files
\colorlet{ink}{plum}
\colorlet{paper}{cream}
\colorlet{mist}{lavender!40!cream}
\definecolor{night}{HTML}{2B2150}
\colorlet{softgray}{plum!60!cream}
% accent of the current part; \partaccent{<colour>} switches it
\colorlet{accent}{butter}
\newcommand{\partaccent}[1]{\colorlet{accent}{#1}}

% ---------------------------------------------------------------------
% Art: \emoji, \mascot (pages of common/art/*.pdf, see build_art.py)
% ---------------------------------------------------------------------
\makeatletter
\newcommand{\artpage}[3]{\expandafter\def\csname idiot@art@#1@#2\endcsname{#3}}
% generated by build_art.py - do not edit
\artpage{emoji}{brain}{1}
\artpage{emoji}{warning}{2}
\artpage{emoji}{rollercoaster}{3}
\artpage{emoji}{doughnut}{4}
\artpage{emoji}{sleeping}{5}
\artpage{emoji}{scream}{6}
\artpage{emoji}{thought}{7}
\artpage{emoji}{question}{8}
\artpage{emoji}{cyclone}{9}
\artpage{emoji}{calendar}{10}
\artpage{emoji}{clock630}{11}
\artpage{emoji}{computer}{12}
\artpage{emoji}{receipt}{13}
\artpage{emoji}{gift}{14}
\artpage{emoji}{lizard}{15}
\artpage{emoji}{boss}{16}
\artpage{emoji}{keyboard}{17}
\artpage{emoji}{wastebasket}{18}
\artpage{emoji}{batterylow}{19}
\artpage{emoji}{skull}{20}
\artpage{emoji}{car}{21}
\artpage{emoji}{ship}{22}
\artpage{emoji}{airplane}{23}
\artpage{emoji}{nauseated}{24}
\artpage{emoji}{vomit}{25}
\artpage{emoji}{eye}{26}
\artpage{emoji}{ear}{27}
\artpage{emoji}{tree}{28}
\artpage{emoji}{wave}{29}
\artpage{emoji}{book}{30}
\artpage{emoji}{walking}{31}
\artpage{emoji}{dizzy}{32}
\artpage{emoji}{cake}{33}
\artpage{emoji}{icecream}{34}
\artpage{emoji}{potato}{35}
\artpage{emoji}{apple}{36}
\artpage{emoji}{burger}{37}
\artpage{emoji}{salad}{38}
\artpage{emoji}{clock12}{39}
\artpage{emoji}{forkknife}{40}
\artpage{emoji}{pasta}{41}
\artpage{emoji}{fire}{42}
\artpage{emoji}{drool}{43}
\artpage{emoji}{pizza}{44}
\artpage{emoji}{bed}{45}
\artpage{emoji}{personbed}{46}
\artpage{emoji}{moon}{47}
\artpage{emoji}{sun}{48}
\artpage{emoji}{zzz}{49}
\artpage{emoji}{dolphin}{50}
\artpage{emoji}{jellyfish}{51}
\artpage{emoji}{rat}{52}
\artpage{emoji}{chipmunk}{53}
\artpage{emoji}{alarm}{54}
\artpage{emoji}{phone}{55}
\artpage{emoji}{ghost}{56}
\artpage{emoji}{alien}{57}
\artpage{emoji}{yawn}{58}
\artpage{emoji}{baby}{59}
\artpage{emoji}{construction}{60}
\artpage{emoji}{broom}{61}
\artpage{emoji}{thermometer}{62}
\artpage{emoji}{cold}{63}
\artpage{emoji}{monkey}{64}
\artpage{emoji}{tiger}{65}
\artpage{emoji}{axe}{66}
\artpage{emoji}{balloon}{67}
\artpage{emoji}{heart}{68}
\artpage{emoji}{anatomicalheart}{69}
\artpage{emoji}{lungs}{70}
\artpage{emoji}{sweat}{71}
\artpage{emoji}{eyes}{72}
\artpage{emoji}{butterfly}{73}
\artpage{emoji}{runner}{74}
\artpage{emoji}{boxing}{75}
\artpage{emoji}{station}{76}
\artpage{emoji}{taxi}{77}
\artpage{emoji}{bell}{78}
\artpage{emoji}{couch}{79}
\artpage{emoji}{muscle}{80}
\artpage{emoji}{toilet}{81}
\artpage{emoji}{floppy}{82}
\artpage{emoji}{door}{83}
\artpage{emoji}{coffee}{84}
\artpage{emoji}{diving}{85}
\artpage{emoji}{wine}{86}
\artpage{emoji}{beer}{87}
\artpage{emoji}{dove}{88}
\artpage{emoji}{puzzle}{89}
\artpage{emoji}{fish}{90}
\artpage{emoji}{trophy}{91}
\artpage{emoji}{detective}{92}
\artpage{emoji}{judge}{93}
\artpage{emoji}{cheese}{94}
\artpage{emoji}{film}{95}
\artpage{emoji}{magnifier}{96}
\artpage{emoji}{key}{97}
\artpage{emoji}{firework}{98}
\artpage{emoji}{house}{99}
\artpage{emoji}{gloves}{100}
\artpage{emoji}{trafficlight}{101}
\artpage{emoji}{stop}{102}
\artpage{emoji}{mute}{103}
\artpage{emoji}{videogame}{104}
\artpage{emoji}{mirror}{105}
\artpage{emoji}{star}{106}
\artpage{emoji}{hospital}{107}
\artpage{emoji}{pill}{108}
\artpage{emoji}{camera}{109}
\artpage{emoji}{tv}{110}
\artpage{emoji}{crown}{111}
\artpage{emoji}{nerd}{112}
\artpage{emoji}{lock}{113}
\artpage{emoji}{musical}{114}
\artpage{emoji}{chips}{115}
\artpage{emoji}{check}{116}
\artpage{emoji}{cross}{117}
\artpage{emoji}{hand}{118}
\artpage{emoji}{wavehand}{119}
\artpage{emoji}{speech}{120}
\artpage{emoji}{pencil}{121}
\artpage{emoji}{memo}{122}
\artpage{emoji}{monocle}{123}
\artpage{emoji}{thinking}{124}
\artpage{emoji}{grin}{125}
\artpage{emoji}{sweatsmile}{126}
\artpage{emoji}{facepalm}{127}
\artpage{emoji}{popcorn}{128}
\artpage{emoji}{bulb}{129}
\artpage{emoji}{sparkles}{130}
\artpage{emoji}{party}{131}
\artpage{emoji}{chart}{132}
\artpage{emoji}{testtube}{133}
\artpage{emoji}{people}{134}
\artpage{emoji}{pin}{135}
\artpage{emoji}{hourglass}{136}
\artpage{emoji}{microscope}{137}
\artpage{emoji}{books}{138}
\artpage{emoji}{megaphone}{139}
\artpage{emoji}{student}{140}
\artpage{emoji}{zany}{141}
\artpage{emoji}{hundred}{142}
\artpage{emoji}{link}{143}
\artpage{emoji}{target}{144}
\artpage{emoji}{clover}{145}
\artpage{emoji}{ufo}{146}
\artpage{emoji}{spider}{147}
\artpage{emoji}{spiderweb}{148}
\artpage{emoji}{snake}{149}
\artpage{emoji}{microphone}{150}
\artpage{emoji}{clown}{151}
\artpage{emoji}{pumpkin}{152}
\artpage{emoji}{knife}{153}
\artpage{emoji}{bat}{154}
\artpage{emoji}{blackcat}{155}
\artpage{emoji}{crossedfingers}{156}
\artpage{emoji}{ladder}{157}
\artpage{emoji}{thumbsup}{158}
\artpage{emoji}{thumbsdown}{159}
\artpage{emoji}{rocket}{160}
\artpage{emoji}{satellite}{161}
\artpage{emoji}{newspaper}{162}
\artpage{emoji}{chocolate}{163}
\artpage{emoji}{syringe}{164}
\artpage{emoji}{dice}{165}
\artpage{emoji}{hauntedhouse}{166}
\artpage{emoji}{flushed}{167}
\artpage{emoji}{fearful}{168}
\artpage{emoji}{graduation}{169}
\artpage{emoji}{teacher}{170}
\artpage{emoji}{scientist}{171}
\artpage{emoji}{chess}{172}
\artpage{emoji}{abacus}{173}
\artpage{emoji}{ruler}{174}
\artpage{emoji}{jeans}{175}
\artpage{emoji}{glasses}{176}
\artpage{emoji}{dna}{177}
\artpage{emoji}{loudspeaker}{178}
\artpage{emoji}{speakinghead}{179}
\artpage{emoji}{seedling}{180}
\artpage{emoji}{weightlifter}{181}
\artpage{emoji}{giraffe}{182}
\artpage{emoji}{mouse}{183}
\artpage{emoji}{exploding}{184}
\artpage{emoji}{owl}{185}
\artpage{emoji}{clipboard}{186}
\artpage{emoji}{numbers}{187}
\artpage{emoji}{medal}{188}
\artpage{emoji}{rose}{189}
\artpage{emoji}{nose}{190}
\artpage{emoji}{tongue}{191}
\artpage{emoji}{yum}{192}
\artpage{emoji}{bread}{193}
\artpage{emoji}{lemon}{194}
\artpage{emoji}{chili}{195}
\artpage{emoji}{salt}{196}
\artpage{emoji}{candy}{197}
\artpage{emoji}{speaker}{198}
\artpage{emoji}{headphones}{199}
\artpage{emoji}{feather}{200}
\artpage{emoji}{gorilla}{201}
\artpage{emoji}{basketball}{202}
\artpage{emoji}{magicwand}{203}
\artpage{emoji}{tophat}{204}
\artpage{emoji}{cloud}{205}
\artpage{emoji}{plug}{206}
\artpage{emoji}{clinking}{207}
\artpage{emoji}{dog}{208}
\artpage{emoji}{rainbow}{209}
\artpage{emoji}{slightsmile}{210}
\artpage{emoji}{palette}{211}
\artpage{emoji}{masks}{212}
\artpage{emoji}{angry}{213}
\artpage{emoji}{rage}{214}
\artpage{emoji}{steam}{215}
\artpage{emoji}{bomb}{216}
\artpage{emoji}{volcano}{217}
\artpage{emoji}{angersymbol}{218}
\artpage{emoji}{crystalball}{219}
\artpage{emoji}{railway}{220}
\artpage{emoji}{worker}{221}
\artpage{emoji}{laughing}{222}
\artpage{emoji}{rofl}{223}
\artpage{emoji}{banana}{224}
\artpage{emoji}{mountain}{225}
\artpage{emoji}{moneybag}{226}
\artpage{emoji}{lightning}{227}
\artpage{emoji}{joker}{228}
\artpage{emoji}{finishflag}{229}
\artpage{emoji}{hug}{230}
\artpage{emoji}{handshake}{231}
\artpage{emoji}{lying}{232}
\artpage{emoji}{neutral}{233}
\artpage{emoji}{carrot}{234}
\artpage{emoji}{wood}{235}
\artpage{emoji}{brokenheart}{236}
\artpage{emoji}{loveletter}{237}
\artpage{emoji}{couple}{238}
\artpage{emoji}{kiss}{239}
\artpage{emoji}{cowboy}{240}
\artpage{emoji}{guitar}{241}
\artpage{emoji}{family}{242}
\artpage{emoji}{police}{243}
\artpage{emoji}{sheep}{244}
\artpage{emoji}{clapping}{245}
\artpage{emoji}{imp}{246}
\artpage{emoji}{ballot}{247}
\artpage{emoji}{smileheart}{248}
\artpage{emoji}{dollar}{249}
\artpage{emoji}{earth}{250}
\artpage{emoji}{raincloud}{251}
\artpage{emoji}{sunbehindcloud}{252}
\artpage{emoji}{pensive}{253}
\artpage{emoji}{cry}{254}
\artpage{emoji}{bandage}{255}
\artpage{emoji}{stethoscope}{256}
\artpage{emoji}{cigarette}{257}
\artpage{emoji}{slotmachine}{258}
\artpage{emoji}{siren}{259}
\artpage{emoji}{sunrise}{260}
\artpage{emoji}{healthworker}{261}
\artpage{emoji}{hearnoevil}{262}
\artpage{emoji}{anchor}{263}
\artpage{emoji}{candle}{264}
\artpage{emoji}{telephone}{265}
\artpage{emoji}{battery}{266}
\artpage{emoji}{weary}{267}
\artpage{emoji}{lotus}{268}
\artpage{mascot}{panic}{1}
\artpage{mascot}{happy}{2}
\artpage{mascot}{scared}{3}
\artpage{mascot}{sleepy}{4}
\artpage{mascot}{sick}{5}
\artpage{mascot}{confused}{6}
\artpage{mascot}{smug}{7}
\artpage{mascot}{hungry}{8}
\artpage{mascot}{thinking}{9}
\artpage{mascot}{sad}{10}
\artpage{mascot}{angry}{11}

\newcommand{\idiot@art}[3]{% kind, name, size options
  \ifcsname idiot@art@#1@#2\endcsname
    \includegraphics[#3, page=\csname idiot@art@#1@#2\endcsname]{#1.pdf}%
  \else
    \PackageError{idiot}{Unknown #1 '#2'; add it to build_art.py and rebuild the art}{}%
  \fi}
% \emoji[<height>]{name}: inline icon, sits on the text baseline like a letter
\newcommand{\emoji}[2][1.15em]{\raisebox{-0.2em}{\idiot@art{emoji}{#2}{height=#1}}}
% \bigemoji[<height>]{name}: no baseline shift, for TikZ nodes and big pictures
\newcommand{\bigemoji}[2][1.5cm]{\idiot@art{emoji}{#2}{height=#1}}
% \mascot[<width>]{mood}: panic, happy, scared, sleepy, sick, confused, smug, hungry, thinking
\newcommand{\mascot}[2][2.6cm]{\idiot@art{mascot}{#2}{width=#1}}
\makeatother
% the older figure files place the mascot with \pic{noggin=<mood>}
\tikzset{pics/noggin/.style={code={\node[inner sep=0, transform shape] {\mascot[2.7cm]{#1}};}}}

% ---------------------------------------------------------------------
% Beamer look
% ---------------------------------------------------------------------
\setbeamertemplate{navigation symbols}{}
\setbeamercolor{normal text}{fg=plum, bg=cream}
\setbeamercolor{structure}{fg=plum}
\setbeamercolor{alerted text}{fg=pinkdark}
\setbeamerfont{alerted text}{series=\bfseries}
\setbeamercolor{frametitle}{fg=plum}
\setbeamersize{text margin left=9mm, text margin right=9mm}
\setbeamertemplate{frame begin}{\raggedright} % justified Armenian leaves big gaps

% cream paper with the banner's faint dot grid
\setbeamertemplate{background}{%
  \begin{tikzpicture}
    \fill[cream] (0,0) rectangle (\paperwidth,\paperheight);
    \fill[pattern={Dots[distance=4.4mm, radius=0.28mm]}, pattern color=plum!9!cream]
      (0,0) rectangle (\paperwidth,\paperheight);
  \end{tikzpicture}}

% topic pill, top right of every frame: set by \topic (emoji + short name)
\def\idiottopicemoji{}
\def\idiottopicname{}
\newcommand{\idiotpill}[3][cardwhite]{% fill, emoji, text
  \tikz[baseline=(p.base)]{\node[fill=#1, draw=plum, line width=0.7pt, rounded corners=2.6mm,
      inner xsep=2.2mm, inner ysep=1.1mm, font=\mediumfont\footnotesize, text=plum] (p)
      {\ifx\relax#2\relax\else\emoji[1.05em]{#2}\hspace{1.2mm}\fi#3};
    \begin{scope}[on background layer]
      \fill[plum, rounded corners=2.6mm] ([xshift=0.7mm, yshift=-0.7mm]p.south west) rectangle ([xshift=0.7mm, yshift=-0.7mm]p.north east);
    \end{scope}}}
% \chip{emoji}{text}: inline sticker label, same as the banner's topic chips
\newcommand{\chip}[3][cardwhite]{\idiotpill[#1]{#2}{#3}}

\setbeamertemplate{frametitle}{%
  \vskip 3.5mm
  \begin{minipage}[t]{0.66\paperwidth}
    \raggedright{\blackfont\fontsize{17.5}{20.5}\selectfont\insertframetitle}\par
    \vskip 1.3mm
    \tikz\fill[accent, rounded corners=0.6mm] (0,0) rectangle (1.5cm,1.1mm);
  \end{minipage}%
  \hfill
  \ifx\idiottopicname\empty\else
    \raisebox{-1mm}{\idiotpill{\idiottopicemoji}{\idiottopicname}}%
  \fi
  \vskip 1mm}

\setbeamertemplate{footline}{%
  \hbox to \paperwidth{\hskip 9mm
    {\mediumfont\fontsize{7}{8}\selectfont\color{plum!55!cream}Ապուշ ուղեղը · \insertshorttitle}\hfill
    {\blackfont\fontsize{8}{9}\selectfont\color{plum!70!cream}\insertframenumber}\hskip 9mm}
  \vskip 3mm}

\setbeamertemplate{itemize item}{\tikz[baseline=-0.65ex]\draw[plum, fill=accent, line width=0.6pt] (0,0) circle (0.95mm);}
\setbeamertemplate{itemize subitem}{\tikz[baseline=-0.65ex]\fill[plum!70] (0,0) circle (0.6mm);}
\setbeamertemplate{enumerate item}{\tikz[baseline=(n.base)]\node[fill=accent, draw=plum, line width=0.6pt,
  circle, inner sep=0.5mm, minimum size=4.2mm, font=\blackfont\scriptsize] (n) {\insertenumlabel};}
\setlength{\leftmargini}{5mm}
\setbeamertemplate{itemize/enumerate body begin}{\setlength{\itemsep}{1.6mm}}
\setbeamertemplate{itemize/enumerate subbody begin}{\setlength{\itemsep}{0.6mm}\small}

% ---------------------------------------------------------------------
% Structure: parts (full-colour openers) and topics (book sections)
% ---------------------------------------------------------------------
% \topic{<emoji>}{<short name>}: starts a book section; the pill shows on every following frame
\newcommand{\topic}[2]{\subsection{#2}\gdef\idiottopicemoji{#1}\gdef\idiottopicname{#2}}
\newcommand{\notopic}{\gdef\idiottopicemoji{}\gdef\idiottopicname{}}

% \partopener{<label>}{<title>}{<subtitle>}{<mood>}{<colour>}
\newcommand{\partopener}[5]{%
  \section{#2}\notopic\partaccent{#5}%
  \begin{frame}[plain]
    \begin{tikzpicture}[remember picture, overlay]
      \fill[#5] (current page.south west) rectangle (current page.north east);
      \fill[pattern={Dots[distance=4.4mm, radius=0.28mm]}, pattern color=plum!12!#5]
        (current page.south west) rectangle (current page.north east);
      \fill[cardwhite, opacity=0.45] ([xshift=-3.3cm, yshift=-0.2cm]current page.east) circle (3.1cm);
      \node[anchor=north west, inner sep=0] at ([xshift=10mm, yshift=-12mm]current page.north west)
        {\idiotpill{}{#1}};
      \node[anchor=west, text=plum, align=left, text width=9.2cm, inner sep=0]
        (t) at ([xshift=10mm, yshift=5mm]current page.west)
        {\titlefont\fontsize{38}{40}\selectfont #2};
      \node[anchor=north west, text=plum, text width=9.2cm, inner sep=0, align=left]
        at ([yshift=-5mm]t.south west) {\mediumfont\large #3};
      \node at ([xshift=-3.3cm, yshift=-0.2cm]current page.east) {\mascot[4.6cm]{#4}};
    \end{tikzpicture}
  \end{frame}}

% \topicopener{<emoji>}{<short name>}{<book's section title>}{<what it is about>}{<mood>}
\newcommand{\topicopener}[5]{%
  \topic{#1}{#2}%
  \begin{frame}[plain]
    \begin{tikzpicture}[remember picture, overlay]
      \fill[accent!55!cream] (current page.south west) rectangle (current page.north east);
      \fill[pattern={Dots[distance=4.4mm, radius=0.28mm]}, pattern color=plum!10!accent!55!cream]
        (current page.south west) rectangle (current page.north east);
      \node[anchor=north west, inner sep=0] at ([xshift=10mm, yshift=-12mm]current page.north west)
        {\idiotpill{#1}{#2}};
      \node[anchor=west, text=plum, text width=9.6cm, align=left, inner sep=0]
        (t) at ([xshift=10mm, yshift=6mm]current page.west)
        {\blackfont\fontsize{24}{28}\selectfont «#3»};
      \node[anchor=north west, text=plum, text width=9.6cm, align=left, inner sep=0]
        at ([yshift=-4mm]t.south west) {\mediumfont\large #4};
      \node at ([xshift=-3.0cm, yshift=-0.4cm]current page.east) {\mascot[4.2cm]{#5}};
      \node[rotate=-12] at ([xshift=-1.4cm, yshift=2.3cm]current page.east) {\bigemoji[1.5cm]{#1}};
    \end{tikzpicture}
  \end{frame}}

% ---------------------------------------------------------------------
% Boxes (sticker style: plum outline + hard plum shadow, like the banner)
% ---------------------------------------------------------------------
\tcbset{
  sticker/.style={enhanced, arc=3mm, boxrule=0.9pt, colframe=plum, colback=cardwhite, halign=left,
    left=3mm, right=3mm, top=2mm, bottom=2mm, fontupper=\normalsize, coltext=plum,
    shadow={1.1mm}{-1.1mm}{0mm}{fill=plum}},
  tagged/.style={sticker, colback=#1!45!cream, colbacktitle=#1, coltitle=plum,
    fonttitle=\blackfont\footnotesize, top=3mm,
    attach boxed title to top left={xshift=4mm, yshift*=-2.4mm},
    boxed title style={arc=2.2mm, boxrule=0.9pt, colframe=plum, left=1.5mm, right=1.5mm, top=0.6mm, bottom=0.6mm}},
}
\newtcolorbox{card}[1][]{sticker, #1}
\newtcolorbox{tryit}[1][]{tagged=mint, title={\emoji[1em]{testtube}\ Փորձիր հիմա}, #1}
\newtcolorbox{poll}[1][]{tagged=butter, title={\emoji[1em]{hand}\ Հարցում}, #1}
\newtcolorbox{chatbox}[1][]{tagged=lavender, title={\emoji[1em]{speech}\ Գրիր չաթում}, #1}
\newtcolorbox{scienceupdate}[1][]{tagged=peach, fontupper=\small, title={\emoji[1em]{microscope}\ Գիտությունն այսօր}, #1}
\newtcolorbox{bookquote}[1][]{sticker, colback=cardwhite, fontupper=\mediumfont, left=4mm,
  borderline west={1.4mm}{0pt}{accent},
  before upper={{\blackfont\footnotesize\emoji[1em]{book}\ Բըրնեթը գրում է}\par\smallskip}, #1}

% Brain.exe-style popup window: \begin{popup}{<window title>} ... \end{popup}
\newcommand{\winbuttons}{%
  \tikz[baseline=-0.6ex, line width=0.6pt]{
    \foreach \i/\f in {0/cardwhite, 1/cardwhite, 2/blush}{
      \filldraw[fill=\f, draw=plum, rounded corners=0.6mm] (\i*4.3mm,-1.6mm) rectangle ++(3.3mm,3.3mm);}
    \draw[plum] (0.8mm,-0.8mm) -- (2.5mm,-0.8mm);
    \draw[plum] (5.1mm,-0.8mm) rectangle (6.8mm,0.9mm);
    \draw[plum] (9.4mm,-0.9mm) -- (11.1mm,0.8mm) (9.4mm,0.8mm) -- (11.1mm,-0.9mm);}}
\newtcolorbox{popup}[2][]{sticker, colback=cardwhite, colbacktitle=brainpink, coltitle=plum,
  fonttitle=\blackfont\small, title={\emoji[1.1em]{warning}\hspace{1.5mm}#2\hfill\winbuttons},
  titlerule=0.9pt, toptitle=0.8mm, bottomtitle=0.8mm, #1}

% Myth or fact: \mythorfact{<statement>}{<verdict>}{<colour>}{<explanation>}; verdict shows on slide 2
\newcommand{\mythorfact}[4]{%
  \begin{center}
  \begin{tikzpicture}
    \node[draw=plum, line width=0.9pt, fill=cardwhite, rounded corners=4mm, inner xsep=8mm, inner ysep=6mm,
          text width=10cm, align=center, font=\blackfont\Large, text=plum] (mf) {«#1»};
    \begin{scope}[on background layer]
      \fill[plum, rounded corners=4mm] ([xshift=1.2mm, yshift=-1.2mm]mf.south west) rectangle ([xshift=1.2mm, yshift=-1.2mm]mf.north east);
    \end{scope}
    \node[anchor=west] at ([xshift=6mm]mf.north west) {\idiotpill[butter]{question}{Առասպել, թե՞ փաստ}};
    \onslide<2->{\node[draw=#3, text=#3, line width=2pt, rounded corners=2mm, rotate=-10, font=\blackfont\LARGE,
      inner sep=4pt, fill=cardwhite] at ([xshift=-8mm, yshift=1mm]mf.south east) {#2};}
  \end{tikzpicture}
  \end{center}
  \onslide<2->{#4}}

% Big question with the confused mascot: \bigquestion{<question>}{<small line>}
\newcommand{\bigquestion}[2]{%
  \vfill
  \begin{center}
  \begin{tikzpicture}
    \node[text width=9.6cm, align=center, font=\blackfont\fontsize{24}{29}\selectfont, text=plum] (q) {#1};
    \node[below=4mm of q, text width=9.6cm, align=center, font=\mediumfont\large, text=softgray] {#2};
    \node[anchor=east] at ([xshift=-2mm]q.west) {\mascot[2.6cm]{confused}};
  \end{tikzpicture}
  \end{center}
  \vfill}

% Small source line at the bottom of a frame: \source{Loftus & Palmer, 1974}
\newcommand{\source}[1]{%
  \begin{tikzpicture}[remember picture, overlay]
    \node[anchor=south east, font=\fontsize{6.5}{7.5}\selectfont, text=plum!60!cream, inner sep=0]
      at ([xshift=-17mm, yshift=3.4mm]current page.south east) {Աղբյուր՝ #1};
  \end{tikzpicture}}

% Speech bubbles (handwritten) and the overlay helper for TikZ:
% \node[visible on=<2->] ... keeps the layout but hides the node until slide 2.
\tikzset{
  say/.style={rectangle callout, draw=plum, line width=0.8pt, fill=cardwhite, rounded corners=3pt,
    inner sep=4pt, align=center, font=\handfont\normalsize, text=plum},
  think/.style={cloud callout, cloud puffs=12, draw=plum, line width=0.8pt, fill=cardwhite,
    inner sep=2pt, align=center, font=\handfont\normalsize, text=plum},
  % boxes for diagrams drawn directly in a frame: a style with #1 cannot be defined inside a frame body
  flow/.style={draw=plum, line width=0.8pt, fill=cardwhite, rounded corners=2mm, align=center,
    font=\small, text=plum, inner sep=4pt, text width=#1},
  flow/.default=3cm,
  chatmsg/.style={draw=plum, line width=0.8pt, rounded corners=3mm, inner sep=6pt, text width=#1,
    align=left, font=\small, text=plum},
  chatmsg/.default=7cm,
  invisible/.style={opacity=0, text opacity=0},
  visible on/.style={alt={#1{}{invisible}}},
  alt/.code args={<#1>#2#3}{\alt<#1>{\pgfkeysalso{#2}}{\pgfkeysalso{#3}}},
}

%  Build from the deck folder:  latexmk -lualatex -outdir=build webinarNN.tex
%  Art (emoji, mascot moods, title banner) comes from common/build_art.py.
%
%  All styling lives here, so decks only hold content. Palette and fonts
%  are the banner's (banner/build_banner.py): plum text on cream, pastel
%  accents, Mardoto for text, Glazz for big titles, Twemoji icons and
%  the pink brain mascot (DECISIONS.md #4-#6).
% =====================================================================

% ---------------------------------------------------------------------
% Fonts (LuaLaTeX + fontspec), shared with the banner
% ---------------------------------------------------------------------
\usepackage{fontspec}
\tracinglostchars=3 % a glyph missing from a font stops the build instead of printing nothing
\defaultfontfeatures{Ligatures=TeX}
\def\idiotfontdir{../../banner/assets/fonts/}
\setsansfont{Mardoto}[Path=\idiotfontdir, Extension=.ttf, UprightFont=*-Regular, BoldFont=*-Bold]
\setmainfont{Mardoto}[Path=\idiotfontdir, Extension=.ttf, UprightFont=*-Regular, BoldFont=*-Bold]
\newfontfamily\mediumfont{Mardoto-Medium}[Path=\idiotfontdir, Extension=.ttf]
\newfontfamily\blackfont{Mardoto-Black}[Path=\idiotfontdir, Extension=.ttf]
\newfontfamily\titlefont{GlazzArm-Regular}[Path=\idiotfontdir, Extension=.ttf] % Armenian letters only
\newfontfamily\handfont{Ink Free} % Windows font with Armenian; speech bubbles only
% Mardoto has no italic: emphasis is medium weight in the accent colour instead.
\renewcommand{\emph}[1]{{\mediumfont\color{plum}#1}}
% No hyphenation: split words are hard to read on a shared screen.
\hyphenpenalty=10000
\exhyphenpenalty=10000

% ---------------------------------------------------------------------
% Packages
% ---------------------------------------------------------------------
\usepackage{graphicx}
\usepackage{booktabs}
\usepackage[most]{tcolorbox}
\usepackage{tikz}
\usetikzlibrary{arrows.meta, positioning, calc, shapes.geometric, shapes.callouts, shapes.misc,
  decorations.pathmorphing, decorations.pathreplacing, backgrounds, fit, patterns.meta}
\usepackage{pgfplots}
\pgfplotsset{compat=1.18}
\usepgfplotslibrary{fillbetween}
\graphicspath{{../common/art/}{images/}}

% ---------------------------------------------------------------------
% Colours: the banner palette. Charts with 3+ series use the Armenian
% flag colours (amRed / amBlue / amOrange).
% ---------------------------------------------------------------------
\definecolor{plum}{HTML}{4A2340}
\definecolor{brainpink}{HTML}{F7A8C4}
\definecolor{fold}{HTML}{E57FA6}
\definecolor{cream}{HTML}{FFF6EC}
\definecolor{cardwhite}{HTML}{FFFDF9}
\definecolor{peach}{HTML}{FFD9C2}
\definecolor{lavender}{HTML}{D9C8F5}
\definecolor{mint}{HTML}{B5E8D5}
\definecolor{butter}{HTML}{FFE8A3}
\definecolor{blush}{HTML}{FF8FA8}
\definecolor{amRed}{HTML}{D90012}
\definecolor{amBlue}{HTML}{0033A0}
\definecolor{amOrange}{HTML}{F2A800}
% darker partners for text on cream
\colorlet{mintdark}{mint!45!black!80!green}
\definecolor{mintdark}{HTML}{2E7D62}
\definecolor{lavenderdark}{HTML}{5B3FA0}
\definecolor{peachdark}{HTML}{B5532A}
\definecolor{pinkdark}{HTML}{C2185B}
% names used by the figure files
\colorlet{ink}{plum}
\colorlet{paper}{cream}
\colorlet{mist}{lavender!40!cream}
\definecolor{night}{HTML}{2B2150}
\colorlet{softgray}{plum!60!cream}
% accent of the current part; \partaccent{<colour>} switches it
\colorlet{accent}{butter}
\newcommand{\partaccent}[1]{\colorlet{accent}{#1}}

% ---------------------------------------------------------------------
% Art: \emoji, \mascot (pages of common/art/*.pdf, see build_art.py)
% ---------------------------------------------------------------------
\makeatletter
\newcommand{\artpage}[3]{\expandafter\def\csname idiot@art@#1@#2\endcsname{#3}}
% generated by build_art.py - do not edit
\artpage{emoji}{brain}{1}
\artpage{emoji}{warning}{2}
\artpage{emoji}{rollercoaster}{3}
\artpage{emoji}{doughnut}{4}
\artpage{emoji}{sleeping}{5}
\artpage{emoji}{scream}{6}
\artpage{emoji}{thought}{7}
\artpage{emoji}{question}{8}
\artpage{emoji}{cyclone}{9}
\artpage{emoji}{calendar}{10}
\artpage{emoji}{clock630}{11}
\artpage{emoji}{computer}{12}
\artpage{emoji}{receipt}{13}
\artpage{emoji}{gift}{14}
\artpage{emoji}{lizard}{15}
\artpage{emoji}{boss}{16}
\artpage{emoji}{keyboard}{17}
\artpage{emoji}{wastebasket}{18}
\artpage{emoji}{batterylow}{19}
\artpage{emoji}{skull}{20}
\artpage{emoji}{car}{21}
\artpage{emoji}{ship}{22}
\artpage{emoji}{airplane}{23}
\artpage{emoji}{nauseated}{24}
\artpage{emoji}{vomit}{25}
\artpage{emoji}{eye}{26}
\artpage{emoji}{ear}{27}
\artpage{emoji}{tree}{28}
\artpage{emoji}{wave}{29}
\artpage{emoji}{book}{30}
\artpage{emoji}{walking}{31}
\artpage{emoji}{dizzy}{32}
\artpage{emoji}{cake}{33}
\artpage{emoji}{icecream}{34}
\artpage{emoji}{potato}{35}
\artpage{emoji}{apple}{36}
\artpage{emoji}{burger}{37}
\artpage{emoji}{salad}{38}
\artpage{emoji}{clock12}{39}
\artpage{emoji}{forkknife}{40}
\artpage{emoji}{pasta}{41}
\artpage{emoji}{fire}{42}
\artpage{emoji}{drool}{43}
\artpage{emoji}{pizza}{44}
\artpage{emoji}{bed}{45}
\artpage{emoji}{personbed}{46}
\artpage{emoji}{moon}{47}
\artpage{emoji}{sun}{48}
\artpage{emoji}{zzz}{49}
\artpage{emoji}{dolphin}{50}
\artpage{emoji}{jellyfish}{51}
\artpage{emoji}{rat}{52}
\artpage{emoji}{chipmunk}{53}
\artpage{emoji}{alarm}{54}
\artpage{emoji}{phone}{55}
\artpage{emoji}{ghost}{56}
\artpage{emoji}{alien}{57}
\artpage{emoji}{yawn}{58}
\artpage{emoji}{baby}{59}
\artpage{emoji}{construction}{60}
\artpage{emoji}{broom}{61}
\artpage{emoji}{thermometer}{62}
\artpage{emoji}{cold}{63}
\artpage{emoji}{monkey}{64}
\artpage{emoji}{tiger}{65}
\artpage{emoji}{axe}{66}
\artpage{emoji}{balloon}{67}
\artpage{emoji}{heart}{68}
\artpage{emoji}{anatomicalheart}{69}
\artpage{emoji}{lungs}{70}
\artpage{emoji}{sweat}{71}
\artpage{emoji}{eyes}{72}
\artpage{emoji}{butterfly}{73}
\artpage{emoji}{runner}{74}
\artpage{emoji}{boxing}{75}
\artpage{emoji}{station}{76}
\artpage{emoji}{taxi}{77}
\artpage{emoji}{bell}{78}
\artpage{emoji}{couch}{79}
\artpage{emoji}{muscle}{80}
\artpage{emoji}{toilet}{81}
\artpage{emoji}{floppy}{82}
\artpage{emoji}{door}{83}
\artpage{emoji}{coffee}{84}
\artpage{emoji}{diving}{85}
\artpage{emoji}{wine}{86}
\artpage{emoji}{beer}{87}
\artpage{emoji}{dove}{88}
\artpage{emoji}{puzzle}{89}
\artpage{emoji}{fish}{90}
\artpage{emoji}{trophy}{91}
\artpage{emoji}{detective}{92}
\artpage{emoji}{judge}{93}
\artpage{emoji}{cheese}{94}
\artpage{emoji}{film}{95}
\artpage{emoji}{magnifier}{96}
\artpage{emoji}{key}{97}
\artpage{emoji}{firework}{98}
\artpage{emoji}{house}{99}
\artpage{emoji}{gloves}{100}
\artpage{emoji}{trafficlight}{101}
\artpage{emoji}{stop}{102}
\artpage{emoji}{mute}{103}
\artpage{emoji}{videogame}{104}
\artpage{emoji}{mirror}{105}
\artpage{emoji}{star}{106}
\artpage{emoji}{hospital}{107}
\artpage{emoji}{pill}{108}
\artpage{emoji}{camera}{109}
\artpage{emoji}{tv}{110}
\artpage{emoji}{crown}{111}
\artpage{emoji}{nerd}{112}
\artpage{emoji}{lock}{113}
\artpage{emoji}{musical}{114}
\artpage{emoji}{chips}{115}
\artpage{emoji}{check}{116}
\artpage{emoji}{cross}{117}
\artpage{emoji}{hand}{118}
\artpage{emoji}{wavehand}{119}
\artpage{emoji}{speech}{120}
\artpage{emoji}{pencil}{121}
\artpage{emoji}{memo}{122}
\artpage{emoji}{monocle}{123}
\artpage{emoji}{thinking}{124}
\artpage{emoji}{grin}{125}
\artpage{emoji}{sweatsmile}{126}
\artpage{emoji}{facepalm}{127}
\artpage{emoji}{popcorn}{128}
\artpage{emoji}{bulb}{129}
\artpage{emoji}{sparkles}{130}
\artpage{emoji}{party}{131}
\artpage{emoji}{chart}{132}
\artpage{emoji}{testtube}{133}
\artpage{emoji}{people}{134}
\artpage{emoji}{pin}{135}
\artpage{emoji}{hourglass}{136}
\artpage{emoji}{microscope}{137}
\artpage{emoji}{books}{138}
\artpage{emoji}{megaphone}{139}
\artpage{emoji}{student}{140}
\artpage{emoji}{zany}{141}
\artpage{emoji}{hundred}{142}
\artpage{emoji}{link}{143}
\artpage{emoji}{target}{144}
\artpage{emoji}{clover}{145}
\artpage{emoji}{ufo}{146}
\artpage{emoji}{spider}{147}
\artpage{emoji}{spiderweb}{148}
\artpage{emoji}{snake}{149}
\artpage{emoji}{microphone}{150}
\artpage{emoji}{clown}{151}
\artpage{emoji}{pumpkin}{152}
\artpage{emoji}{knife}{153}
\artpage{emoji}{bat}{154}
\artpage{emoji}{blackcat}{155}
\artpage{emoji}{crossedfingers}{156}
\artpage{emoji}{ladder}{157}
\artpage{emoji}{thumbsup}{158}
\artpage{emoji}{thumbsdown}{159}
\artpage{emoji}{rocket}{160}
\artpage{emoji}{satellite}{161}
\artpage{emoji}{newspaper}{162}
\artpage{emoji}{chocolate}{163}
\artpage{emoji}{syringe}{164}
\artpage{emoji}{dice}{165}
\artpage{emoji}{hauntedhouse}{166}
\artpage{emoji}{flushed}{167}
\artpage{emoji}{fearful}{168}
\artpage{emoji}{graduation}{169}
\artpage{emoji}{teacher}{170}
\artpage{emoji}{scientist}{171}
\artpage{emoji}{chess}{172}
\artpage{emoji}{abacus}{173}
\artpage{emoji}{ruler}{174}
\artpage{emoji}{jeans}{175}
\artpage{emoji}{glasses}{176}
\artpage{emoji}{dna}{177}
\artpage{emoji}{loudspeaker}{178}
\artpage{emoji}{speakinghead}{179}
\artpage{emoji}{seedling}{180}
\artpage{emoji}{weightlifter}{181}
\artpage{emoji}{giraffe}{182}
\artpage{emoji}{mouse}{183}
\artpage{emoji}{exploding}{184}
\artpage{emoji}{owl}{185}
\artpage{emoji}{clipboard}{186}
\artpage{emoji}{numbers}{187}
\artpage{emoji}{medal}{188}
\artpage{emoji}{rose}{189}
\artpage{emoji}{nose}{190}
\artpage{emoji}{tongue}{191}
\artpage{emoji}{yum}{192}
\artpage{emoji}{bread}{193}
\artpage{emoji}{lemon}{194}
\artpage{emoji}{chili}{195}
\artpage{emoji}{salt}{196}
\artpage{emoji}{candy}{197}
\artpage{emoji}{speaker}{198}
\artpage{emoji}{headphones}{199}
\artpage{emoji}{feather}{200}
\artpage{emoji}{gorilla}{201}
\artpage{emoji}{basketball}{202}
\artpage{emoji}{magicwand}{203}
\artpage{emoji}{tophat}{204}
\artpage{emoji}{cloud}{205}
\artpage{emoji}{plug}{206}
\artpage{emoji}{clinking}{207}
\artpage{emoji}{dog}{208}
\artpage{emoji}{rainbow}{209}
\artpage{emoji}{slightsmile}{210}
\artpage{emoji}{palette}{211}
\artpage{emoji}{masks}{212}
\artpage{emoji}{angry}{213}
\artpage{emoji}{rage}{214}
\artpage{emoji}{steam}{215}
\artpage{emoji}{bomb}{216}
\artpage{emoji}{volcano}{217}
\artpage{emoji}{angersymbol}{218}
\artpage{emoji}{crystalball}{219}
\artpage{emoji}{railway}{220}
\artpage{emoji}{worker}{221}
\artpage{emoji}{laughing}{222}
\artpage{emoji}{rofl}{223}
\artpage{emoji}{banana}{224}
\artpage{emoji}{mountain}{225}
\artpage{emoji}{moneybag}{226}
\artpage{emoji}{lightning}{227}
\artpage{emoji}{joker}{228}
\artpage{emoji}{finishflag}{229}
\artpage{emoji}{hug}{230}
\artpage{emoji}{handshake}{231}
\artpage{emoji}{lying}{232}
\artpage{emoji}{neutral}{233}
\artpage{emoji}{carrot}{234}
\artpage{emoji}{wood}{235}
\artpage{emoji}{brokenheart}{236}
\artpage{emoji}{loveletter}{237}
\artpage{emoji}{couple}{238}
\artpage{emoji}{kiss}{239}
\artpage{emoji}{cowboy}{240}
\artpage{emoji}{guitar}{241}
\artpage{emoji}{family}{242}
\artpage{emoji}{police}{243}
\artpage{emoji}{sheep}{244}
\artpage{emoji}{clapping}{245}
\artpage{emoji}{imp}{246}
\artpage{emoji}{ballot}{247}
\artpage{emoji}{smileheart}{248}
\artpage{emoji}{dollar}{249}
\artpage{emoji}{earth}{250}
\artpage{emoji}{raincloud}{251}
\artpage{emoji}{sunbehindcloud}{252}
\artpage{emoji}{pensive}{253}
\artpage{emoji}{cry}{254}
\artpage{emoji}{bandage}{255}
\artpage{emoji}{stethoscope}{256}
\artpage{emoji}{cigarette}{257}
\artpage{emoji}{slotmachine}{258}
\artpage{emoji}{siren}{259}
\artpage{emoji}{sunrise}{260}
\artpage{emoji}{healthworker}{261}
\artpage{emoji}{hearnoevil}{262}
\artpage{emoji}{anchor}{263}
\artpage{emoji}{candle}{264}
\artpage{emoji}{telephone}{265}
\artpage{emoji}{battery}{266}
\artpage{emoji}{weary}{267}
\artpage{emoji}{lotus}{268}
\artpage{mascot}{panic}{1}
\artpage{mascot}{happy}{2}
\artpage{mascot}{scared}{3}
\artpage{mascot}{sleepy}{4}
\artpage{mascot}{sick}{5}
\artpage{mascot}{confused}{6}
\artpage{mascot}{smug}{7}
\artpage{mascot}{hungry}{8}
\artpage{mascot}{thinking}{9}
\artpage{mascot}{sad}{10}
\artpage{mascot}{angry}{11}

\newcommand{\idiot@art}[3]{% kind, name, size options
  \ifcsname idiot@art@#1@#2\endcsname
    \includegraphics[#3, page=\csname idiot@art@#1@#2\endcsname]{#1.pdf}%
  \else
    \PackageError{idiot}{Unknown #1 '#2'; add it to build_art.py and rebuild the art}{}%
  \fi}
% \emoji[<height>]{name}: inline icon, sits on the text baseline like a letter
\newcommand{\emoji}[2][1.15em]{\raisebox{-0.2em}{\idiot@art{emoji}{#2}{height=#1}}}
% \bigemoji[<height>]{name}: no baseline shift, for TikZ nodes and big pictures
\newcommand{\bigemoji}[2][1.5cm]{\idiot@art{emoji}{#2}{height=#1}}
% \mascot[<width>]{mood}: panic, happy, scared, sleepy, sick, confused, smug, hungry, thinking
\newcommand{\mascot}[2][2.6cm]{\idiot@art{mascot}{#2}{width=#1}}
\makeatother
% the older figure files place the mascot with \pic{noggin=<mood>}
\tikzset{pics/noggin/.style={code={\node[inner sep=0, transform shape] {\mascot[2.7cm]{#1}};}}}

% ---------------------------------------------------------------------
% Beamer look
% ---------------------------------------------------------------------
\setbeamertemplate{navigation symbols}{}
\setbeamercolor{normal text}{fg=plum, bg=cream}
\setbeamercolor{structure}{fg=plum}
\setbeamercolor{alerted text}{fg=pinkdark}
\setbeamerfont{alerted text}{series=\bfseries}
\setbeamercolor{frametitle}{fg=plum}
\setbeamersize{text margin left=9mm, text margin right=9mm}
\setbeamertemplate{frame begin}{\raggedright} % justified Armenian leaves big gaps

% cream paper with the banner's faint dot grid
\setbeamertemplate{background}{%
  \begin{tikzpicture}
    \fill[cream] (0,0) rectangle (\paperwidth,\paperheight);
    \fill[pattern={Dots[distance=4.4mm, radius=0.28mm]}, pattern color=plum!9!cream]
      (0,0) rectangle (\paperwidth,\paperheight);
  \end{tikzpicture}}

% topic pill, top right of every frame: set by \topic (emoji + short name)
\def\idiottopicemoji{}
\def\idiottopicname{}
\newcommand{\idiotpill}[3][cardwhite]{% fill, emoji, text
  \tikz[baseline=(p.base)]{\node[fill=#1, draw=plum, line width=0.7pt, rounded corners=2.6mm,
      inner xsep=2.2mm, inner ysep=1.1mm, font=\mediumfont\footnotesize, text=plum] (p)
      {\ifx\relax#2\relax\else\emoji[1.05em]{#2}\hspace{1.2mm}\fi#3};
    \begin{scope}[on background layer]
      \fill[plum, rounded corners=2.6mm] ([xshift=0.7mm, yshift=-0.7mm]p.south west) rectangle ([xshift=0.7mm, yshift=-0.7mm]p.north east);
    \end{scope}}}
% \chip{emoji}{text}: inline sticker label, same as the banner's topic chips
\newcommand{\chip}[3][cardwhite]{\idiotpill[#1]{#2}{#3}}

\setbeamertemplate{frametitle}{%
  \vskip 3.5mm
  \begin{minipage}[t]{0.66\paperwidth}
    \raggedright{\blackfont\fontsize{17.5}{20.5}\selectfont\insertframetitle}\par
    \vskip 1.3mm
    \tikz\fill[accent, rounded corners=0.6mm] (0,0) rectangle (1.5cm,1.1mm);
  \end{minipage}%
  \hfill
  \ifx\idiottopicname\empty\else
    \raisebox{-1mm}{\idiotpill{\idiottopicemoji}{\idiottopicname}}%
  \fi
  \vskip 1mm}

\setbeamertemplate{footline}{%
  \hbox to \paperwidth{\hskip 9mm
    {\mediumfont\fontsize{7}{8}\selectfont\color{plum!55!cream}Ապուշ ուղեղը · \insertshorttitle}\hfill
    {\blackfont\fontsize{8}{9}\selectfont\color{plum!70!cream}\insertframenumber}\hskip 9mm}
  \vskip 3mm}

\setbeamertemplate{itemize item}{\tikz[baseline=-0.65ex]\draw[plum, fill=accent, line width=0.6pt] (0,0) circle (0.95mm);}
\setbeamertemplate{itemize subitem}{\tikz[baseline=-0.65ex]\fill[plum!70] (0,0) circle (0.6mm);}
\setbeamertemplate{enumerate item}{\tikz[baseline=(n.base)]\node[fill=accent, draw=plum, line width=0.6pt,
  circle, inner sep=0.5mm, minimum size=4.2mm, font=\blackfont\scriptsize] (n) {\insertenumlabel};}
\setlength{\leftmargini}{5mm}
\setbeamertemplate{itemize/enumerate body begin}{\setlength{\itemsep}{1.6mm}}
\setbeamertemplate{itemize/enumerate subbody begin}{\setlength{\itemsep}{0.6mm}\small}

% ---------------------------------------------------------------------
% Structure: parts (full-colour openers) and topics (book sections)
% ---------------------------------------------------------------------
% \topic{<emoji>}{<short name>}: starts a book section; the pill shows on every following frame
\newcommand{\topic}[2]{\subsection{#2}\gdef\idiottopicemoji{#1}\gdef\idiottopicname{#2}}
\newcommand{\notopic}{\gdef\idiottopicemoji{}\gdef\idiottopicname{}}

% \partopener{<label>}{<title>}{<subtitle>}{<mood>}{<colour>}
\newcommand{\partopener}[5]{%
  \section{#2}\notopic\partaccent{#5}%
  \begin{frame}[plain]
    \begin{tikzpicture}[remember picture, overlay]
      \fill[#5] (current page.south west) rectangle (current page.north east);
      \fill[pattern={Dots[distance=4.4mm, radius=0.28mm]}, pattern color=plum!12!#5]
        (current page.south west) rectangle (current page.north east);
      \fill[cardwhite, opacity=0.45] ([xshift=-3.3cm, yshift=-0.2cm]current page.east) circle (3.1cm);
      \node[anchor=north west, inner sep=0] at ([xshift=10mm, yshift=-12mm]current page.north west)
        {\idiotpill{}{#1}};
      \node[anchor=west, text=plum, align=left, text width=9.2cm, inner sep=0]
        (t) at ([xshift=10mm, yshift=5mm]current page.west)
        {\titlefont\fontsize{38}{40}\selectfont #2};
      \node[anchor=north west, text=plum, text width=9.2cm, inner sep=0, align=left]
        at ([yshift=-5mm]t.south west) {\mediumfont\large #3};
      \node at ([xshift=-3.3cm, yshift=-0.2cm]current page.east) {\mascot[4.6cm]{#4}};
    \end{tikzpicture}
  \end{frame}}

% \topicopener{<emoji>}{<short name>}{<book's section title>}{<what it is about>}{<mood>}
\newcommand{\topicopener}[5]{%
  \topic{#1}{#2}%
  \begin{frame}[plain]
    \begin{tikzpicture}[remember picture, overlay]
      \fill[accent!55!cream] (current page.south west) rectangle (current page.north east);
      \fill[pattern={Dots[distance=4.4mm, radius=0.28mm]}, pattern color=plum!10!accent!55!cream]
        (current page.south west) rectangle (current page.north east);
      \node[anchor=north west, inner sep=0] at ([xshift=10mm, yshift=-12mm]current page.north west)
        {\idiotpill{#1}{#2}};
      \node[anchor=west, text=plum, text width=9.6cm, align=left, inner sep=0]
        (t) at ([xshift=10mm, yshift=6mm]current page.west)
        {\blackfont\fontsize{24}{28}\selectfont «#3»};
      \node[anchor=north west, text=plum, text width=9.6cm, align=left, inner sep=0]
        at ([yshift=-4mm]t.south west) {\mediumfont\large #4};
      \node at ([xshift=-3.0cm, yshift=-0.4cm]current page.east) {\mascot[4.2cm]{#5}};
      \node[rotate=-12] at ([xshift=-1.4cm, yshift=2.3cm]current page.east) {\bigemoji[1.5cm]{#1}};
    \end{tikzpicture}
  \end{frame}}

% ---------------------------------------------------------------------
% Boxes (sticker style: plum outline + hard plum shadow, like the banner)
% ---------------------------------------------------------------------
\tcbset{
  sticker/.style={enhanced, arc=3mm, boxrule=0.9pt, colframe=plum, colback=cardwhite, halign=left,
    left=3mm, right=3mm, top=2mm, bottom=2mm, fontupper=\normalsize, coltext=plum,
    shadow={1.1mm}{-1.1mm}{0mm}{fill=plum}},
  tagged/.style={sticker, colback=#1!45!cream, colbacktitle=#1, coltitle=plum,
    fonttitle=\blackfont\footnotesize, top=3mm,
    attach boxed title to top left={xshift=4mm, yshift*=-2.4mm},
    boxed title style={arc=2.2mm, boxrule=0.9pt, colframe=plum, left=1.5mm, right=1.5mm, top=0.6mm, bottom=0.6mm}},
}
\newtcolorbox{card}[1][]{sticker, #1}
\newtcolorbox{tryit}[1][]{tagged=mint, title={\emoji[1em]{testtube}\ Փորձիր հիմա}, #1}
\newtcolorbox{poll}[1][]{tagged=butter, title={\emoji[1em]{hand}\ Հարցում}, #1}
\newtcolorbox{chatbox}[1][]{tagged=lavender, title={\emoji[1em]{speech}\ Գրիր չաթում}, #1}
\newtcolorbox{scienceupdate}[1][]{tagged=peach, fontupper=\small, title={\emoji[1em]{microscope}\ Գիտությունն այսօր}, #1}
\newtcolorbox{bookquote}[1][]{sticker, colback=cardwhite, fontupper=\mediumfont, left=4mm,
  borderline west={1.4mm}{0pt}{accent},
  before upper={{\blackfont\footnotesize\emoji[1em]{book}\ Բըրնեթը գրում է}\par\smallskip}, #1}

% Brain.exe-style popup window: \begin{popup}{<window title>} ... \end{popup}
\newcommand{\winbuttons}{%
  \tikz[baseline=-0.6ex, line width=0.6pt]{
    \foreach \i/\f in {0/cardwhite, 1/cardwhite, 2/blush}{
      \filldraw[fill=\f, draw=plum, rounded corners=0.6mm] (\i*4.3mm,-1.6mm) rectangle ++(3.3mm,3.3mm);}
    \draw[plum] (0.8mm,-0.8mm) -- (2.5mm,-0.8mm);
    \draw[plum] (5.1mm,-0.8mm) rectangle (6.8mm,0.9mm);
    \draw[plum] (9.4mm,-0.9mm) -- (11.1mm,0.8mm) (9.4mm,0.8mm) -- (11.1mm,-0.9mm);}}
\newtcolorbox{popup}[2][]{sticker, colback=cardwhite, colbacktitle=brainpink, coltitle=plum,
  fonttitle=\blackfont\small, title={\emoji[1.1em]{warning}\hspace{1.5mm}#2\hfill\winbuttons},
  titlerule=0.9pt, toptitle=0.8mm, bottomtitle=0.8mm, #1}

% Myth or fact: \mythorfact{<statement>}{<verdict>}{<colour>}{<explanation>}; verdict shows on slide 2
\newcommand{\mythorfact}[4]{%
  \begin{center}
  \begin{tikzpicture}
    \node[draw=plum, line width=0.9pt, fill=cardwhite, rounded corners=4mm, inner xsep=8mm, inner ysep=6mm,
          text width=10cm, align=center, font=\blackfont\Large, text=plum] (mf) {«#1»};
    \begin{scope}[on background layer]
      \fill[plum, rounded corners=4mm] ([xshift=1.2mm, yshift=-1.2mm]mf.south west) rectangle ([xshift=1.2mm, yshift=-1.2mm]mf.north east);
    \end{scope}
    \node[anchor=west] at ([xshift=6mm]mf.north west) {\idiotpill[butter]{question}{Առասպել, թե՞ փաստ}};
    \onslide<2->{\node[draw=#3, text=#3, line width=2pt, rounded corners=2mm, rotate=-10, font=\blackfont\LARGE,
      inner sep=4pt, fill=cardwhite] at ([xshift=-8mm, yshift=1mm]mf.south east) {#2};}
  \end{tikzpicture}
  \end{center}
  \onslide<2->{#4}}

% Big question with the confused mascot: \bigquestion{<question>}{<small line>}
\newcommand{\bigquestion}[2]{%
  \vfill
  \begin{center}
  \begin{tikzpicture}
    \node[text width=9.6cm, align=center, font=\blackfont\fontsize{24}{29}\selectfont, text=plum] (q) {#1};
    \node[below=4mm of q, text width=9.6cm, align=center, font=\mediumfont\large, text=softgray] {#2};
    \node[anchor=east] at ([xshift=-2mm]q.west) {\mascot[2.6cm]{confused}};
  \end{tikzpicture}
  \end{center}
  \vfill}

% Small source line at the bottom of a frame: \source{Loftus & Palmer, 1974}
\newcommand{\source}[1]{%
  \begin{tikzpicture}[remember picture, overlay]
    \node[anchor=south east, font=\fontsize{6.5}{7.5}\selectfont, text=plum!60!cream, inner sep=0]
      at ([xshift=-17mm, yshift=3.4mm]current page.south east) {Աղբյուր՝ #1};
  \end{tikzpicture}}

% Speech bubbles (handwritten) and the overlay helper for TikZ:
% \node[visible on=<2->] ... keeps the layout but hides the node until slide 2.
\tikzset{
  say/.style={rectangle callout, draw=plum, line width=0.8pt, fill=cardwhite, rounded corners=3pt,
    inner sep=4pt, align=center, font=\handfont\normalsize, text=plum},
  think/.style={cloud callout, cloud puffs=12, draw=plum, line width=0.8pt, fill=cardwhite,
    inner sep=2pt, align=center, font=\handfont\normalsize, text=plum},
  % boxes for diagrams drawn directly in a frame: a style with #1 cannot be defined inside a frame body
  flow/.style={draw=plum, line width=0.8pt, fill=cardwhite, rounded corners=2mm, align=center,
    font=\small, text=plum, inner sep=4pt, text width=#1},
  flow/.default=3cm,
  chatmsg/.style={draw=plum, line width=0.8pt, rounded corners=3mm, inner sep=6pt, text width=#1,
    align=left, font=\small, text=plum},
  chatmsg/.default=7cm,
  invisible/.style={opacity=0, text opacity=0},
  visible on/.style={alt={#1{}{invisible}}},
  alt/.code args={<#1>#2#3}{\alt<#1>{\pgfkeysalso{#2}}{\pgfkeysalso{#3}}},
}

%  Build with LuaLaTeX:  latexmk -lualatex chNN.tex
%
%  ALL styling lives here (colours, fonts, theme tweaks, mascot,
%  reusable boxes). Chapter files should only contain content, so the
%  whole series can be restyled from this one file.
% =====================================================================

% ---------------------------------------------------------------------
% Theme
% metropolis is in TeX Live but no longer maintained; "moloch" (also in
% TeX Live 2025) is its maintained fork and should be a drop-in swap if
% metropolis ever breaks:  \usetheme{moloch}
% ---------------------------------------------------------------------
\usetheme[
  progressbar=frametitle,
  numbering=counter,
  sectionpage=none,      % we draw our own section openers (\sectionopener)
  block=fill
]{metropolis}

% ---------------------------------------------------------------------
% Fonts (LuaLaTeX + fontspec)
% Segoe UI is the body font: clean on screen and it ships with Windows
% WITH Armenian glyphs (checked with fontTools, U+0531-U+0587 present in
% regular/bold/italic/semibold). So the deck can be translated to
% Armenian later without touching fonts. Fira Sans (metropolis' default)
% and DejaVu/Noto are NOT installed on this machine.
% Ink Free (handwritten look, also has Armenian) is used for the mascot's
% speech bubbles and playful labels.
% If you build on a non-Windows machine: swap Segoe UI for
% "Noto Sans" + "Noto Sans Armenian" (luaotfload fallback) or DejaVu Sans.
% ---------------------------------------------------------------------
\usepackage{fontspec}
\setsansfont{Segoe UI}[
  BoldFont       = {Segoe UI Bold},
  ItalicFont     = {Segoe UI Italic},
  BoldItalicFont = {Segoe UI Bold Italic}
]
\setmainfont{Segoe UI}[
  BoldFont       = {Segoe UI Bold},
  ItalicFont     = {Segoe UI Italic},
  BoldItalicFont = {Segoe UI Bold Italic}
]
\setmonofont{Consolas}
% \newfontfamily does not turn on TeX ligatures by default, so `` '' would
% print literally without Ligatures=TeX.
\newfontfamily\handfont{Ink Free}[Ligatures=TeX, BoldFont={Ink Free}]
\newfontfamily\semiboldfont{Segoe UI Semibold}[Ligatures=TeX]

% No hyphenation on slides: split words are hard to read on a shared screen.
\hyphenpenalty=10000
\exhyphenpenalty=10000

% ---------------------------------------------------------------------
% Packages
% ---------------------------------------------------------------------
\usepackage{graphicx}
\usepackage{booktabs}
\usepackage{tabularx}
\usepackage[most]{tcolorbox}

% ---------------------------------------------------------------------
% Icons (Font Awesome 5 Free Solid)
% The fontawesome5 PACKAGE fails under LuaLaTeX on this minimal TeX Live
% (it needs luatexbase.sty, which is not installed). So we load the OTF
% that ships with it directly and define the icons we use by code point.
% Names match the fontawesome5 package, so switching back is a one-liner.
% Code points were read from the font's cmap with fontTools.
% To add one: find the glyph name on fontawesome.com/v5 and its code point.
% ---------------------------------------------------------------------
\newfontfamily\FAsolid{FontAwesome5Free-Solid-900}[Extension=.otf,
  BoldFont=FontAwesome5Free-Solid-900]
\newcommand{\faicon}[1]{{\FAsolid\symbol{"#1}}}
\newcommand{\faBrain}{\faicon{F5DC}}
\newcommand{\faCarSide}{\faicon{F5E4}}
\newcommand{\faCar}{\faicon{F1B9}}
\newcommand{\faBus}{\faicon{F207}}
\newcommand{\faShip}{\faicon{F21A}}
\newcommand{\faPlane}{\faicon{F072}}
\newcommand{\faVrCardboard}{\faicon{F729}}
\newcommand{\faHamburger}{\faicon{F805}}
\newcommand{\faIceCream}{\faicon{F810}}
\newcommand{\faCookieBite}{\faicon{F564}}
\newcommand{\faUtensils}{\faicon{F2E7}}
\newcommand{\faPizzaSlice}{\faicon{F818}}
\newcommand{\faAppleAlt}{\faicon{F5D1}}
\newcommand{\faCarrot}{\faicon{F787}}
\newcommand{\faCandyCane}{\faicon{F786}}
\newcommand{\faBed}{\faicon{F236}}
\newcommand{\faMoon}{\faicon{F186}}
\newcommand{\faSun}{\faicon{F185}}
\newcommand{\faRunning}{\faicon{F70C}}
\newcommand{\faWalking}{\faicon{F554}}
\newcommand{\faEye}{\faicon{F06E}}
\newcommand{\faEyeSlash}{\faicon{F070}}
\newcommand{\faBolt}{\faicon{F0E7}}
\newcommand{\faHeartbeat}{\faicon{F21E}}
\newcommand{\faHeart}{\faicon{F004}}
\newcommand{\faLungs}{\faicon{F604}}
\newcommand{\faTint}{\faicon{F043}}
\newcommand{\faHeadphones}{\faicon{F025}}
\newcommand{\faDeaf}{\faicon{F2A4}}
\newcommand{\faClock}{\faicon{F017}}
\newcommand{\faHourglassHalf}{\faicon{F252}}
\newcommand{\faCommentDots}{\faicon{F4AD}}
\newcommand{\faPoll}{\faicon{F681}}
\newcommand{\faFlask}{\faicon{F0C3}}
\newcommand{\faQuestionCircle}{\faicon{F059}}
\newcommand{\faSyncAlt}{\faicon{F2F1}}
\newcommand{\faGhost}{\faicon{F6E2}}
\newcommand{\faSkull}{\faicon{F54C}}
\newcommand{\faCat}{\faicon{F6BE}}
\newcommand{\faSmile}{\faicon{F118}}
\newcommand{\faLaugh}{\faicon{F599}}
\newcommand{\faDizzy}{\faicon{F567}}
\newcommand{\faTired}{\faicon{F5C8}}
\newcommand{\faSurprise}{\faicon{F5C2}}
\newcommand{\faGrimace}{\faicon{F57F}}
\newcommand{\faFrown}{\faicon{F119}}
\newcommand{\faSadTear}{\faicon{F5B4}}
\newcommand{\faAngry}{\faicon{F556}}
\newcommand{\faTaxi}{\faicon{F1BA}}
\newcommand{\faSubway}{\faicon{F239}}
\newcommand{\faBell}{\faicon{F0F3}}
\newcommand{\faSearch}{\faicon{F002}}
\newcommand{\faChild}{\faicon{F1AE}}
\newcommand{\faBaby}{\faicon{F77C}}
\newcommand{\faUser}{\faicon{F007}}
\newcommand{\faBook}{\faicon{F02D}}
\newcommand{\faMobileAlt}{\faicon{F3CD}}
\newcommand{\faCouch}{\faicon{F4B8}}
\newcommand{\faExclamationTriangle}{\faicon{F071}}
\newcommand{\faCheck}{\faicon{F00C}}
\newcommand{\faTimes}{\faicon{F00D}}
\newcommand{\faHome}{\faicon{F015}}
\newcommand{\faUserSecret}{\faicon{F21B}}
\newcommand{\faDoorOpen}{\faicon{F52B}}
\newcommand{\faLightbulb}{\faicon{F0EB}}
\newcommand{\faFish}{\faicon{F578}}
\newcommand{\faWater}{\faicon{F773}}
\newcommand{\faTree}{\faicon{F1BB}}
\newcommand{\faHandPaper}{\faicon{F256}}
\newcommand{\faCalendarAlt}{\faicon{F073}}
\newcommand{\faMicroscope}{\faicon{F610}}
\newcommand{\faDna}{\faicon{F471}}
\newcommand{\faArrowRight}{\faicon{F061}}
\newcommand{\faRedo}{\faicon{F01E}}
\newcommand{\faFire}{\faicon{F06D}}
\newcommand{\faWind}{\faicon{F72E}}
\newcommand{\faGlasses}{\faicon{F530}}
\newcommand{\faTheaterMasks}{\faicon{F630}}
\newcommand{\faAward}{\faicon{F559}}
\newcommand{\faBullseye}{\faicon{F140}}
\newcommand{\faGrinTongue}{\faicon{F589}}
\newcommand{\faMehRollingEyes}{\faicon{F5A5}}
\usepackage{tikz}
\usetikzlibrary{
  arrows.meta, positioning, calc, shapes.geometric, shapes.callouts,
  shapes.symbols, shapes.misc, decorations.pathmorphing,
  decorations.pathreplacing, backgrounds, fit, shadows.blur
}
\usepackage{pgfplots}
\pgfplotsset{compat=1.18}
\usepgfplotslibrary{fillbetween}

% ---------------------------------------------------------------------
% Colours
% Armenian flag colours are the accent palette (also used for any chart
% with 3+ colours): red #D90012, blue #0033A0, orange #F2A800.
% ---------------------------------------------------------------------
\definecolor{amRed}{HTML}{D90012}
\definecolor{amBlue}{HTML}{0033A0}
\definecolor{amOrange}{HTML}{F2A800}
\definecolor{ink}{HTML}{1B1B2F}       % near-black navy for text
\definecolor{paper}{HTML}{FFFBF2}     % warm off-white background
\definecolor{mist}{HTML}{ECEFF6}      % light grey-blue panels
\definecolor{night}{HTML}{18224F}     % sleep section
\definecolor{softgray}{HTML}{8A8FA3}
\definecolor{brainpink}{HTML}{F7B5C8} % mascot body
\definecolor{brainline}{HTML}{B03A62} % mascot outline
\definecolor{sickgreen}{HTML}{BFE0A8}
\definecolor{sickline}{HTML}{5E8C46}

% Beamer colour slots
\setbeamercolor{normal text}{fg=ink, bg=paper}
\setbeamercolor{background canvas}{bg=paper}
\setbeamercolor{alerted text}{fg=amRed}
\setbeamercolor{example text}{fg=amBlue}
\setbeamercolor{frametitle}{fg=white, bg=ink}
\setbeamercolor{progress bar}{fg=amOrange, bg=amOrange!25}
\setbeamercolor{title separator}{fg=amOrange}
\setbeamercolor{progress bar in head/foot}{fg=amOrange, bg=amOrange!25}
\setbeamercolor{progress bar in section page}{fg=amOrange, bg=amOrange!25}
\setbeamercolor{page number in head/foot}{fg=softgray}
\setbeamercolor{block title}{fg=white, bg=amBlue}
\setbeamercolor{block body}{fg=ink, bg=mist}
\setbeamercolor{itemize item}{fg=amOrange}
\setbeamercolor{itemize subitem}{fg=amOrange}

% Font sizes / weights
\setbeamerfont{frametitle}{size=\Large}
\setbeamerfont{title}{size=\huge, series=\bfseries}
\setbeamerfont{subtitle}{size=\large}
\setbeamerfont{page number in head/foot}{size=\scriptsize}
\setbeamerfont{alerted text}{series=\bfseries}
\setbeamertemplate{itemize item}{\raisebox{0.1ex}{\scriptsize$\blacktriangleright$}}
\setbeamertemplate{itemize subitem}{\raisebox{0.1ex}{\tiny$\bullet$}}
\setbeamertemplate{navigation symbols}{}

% Per-section accent: changes the frame-title bar from this point on.
% \sectioncolors{<bar colour>}{<bar text colour>}
\newcommand{\sectioncolors}[2]{%
  \setbeamercolor{frametitle}{fg=#2, bg=#1}%
  \colorlet{sectionaccent}{#1}%
}
\colorlet{sectionaccent}{ink}

% Handy inline helpers
\newcommand{\hl}[1]{\textcolor{amRed}{\textbf{#1}}}       % red highlight
\newcommand{\hlb}[1]{\textcolor{amBlue}{\textbf{#1}}}     % blue highlight
\newcommand{\hlo}[1]{\textcolor{amOrange!85!black}{\textbf{#1}}} % orange
\newcommand{\hand}[1]{{\handfont #1}}                      % handwritten
\newcommand{\muted}[1]{\textcolor{softgray}{#1}}

% =====================================================================
% Mascot: "Noggin", an original cartoon brain (own design: puffy pink
% cloud-outline brain, big eyes, little brainstem "neck").
%   \noggin[<scale>]{<mood>}     inline tikzpicture
%   \pic[scale=..] {noggin=<mood>};   inside another tikzpicture
% moods: happy, panic, sleepy, confused, sick, hungry, smug
% =====================================================================
\makeatletter
\newcommand{\noggin@body}[2]{% #1 fill, #2 line
  % brainstem / neck (drawn first so the body overlaps it)
  \filldraw[fill=#1!80!brainline, draw=#2, line width=0.9pt, rounded corners=2pt]
    (-0.17,-0.55) -- (-0.13,-1.12) -- (0.13,-1.12) -- (0.17,-0.55) -- cycle;
  % cerebellum bump
  \filldraw[fill=#1!85!brainline, draw=#2, line width=0.9pt]
    (0.62,-0.55) ellipse (0.36 and 0.24);
  % puffy cloud-like outline
  \filldraw[fill=#1, draw=#2, line width=1.1pt]
    plot[domain=0:360, samples=180, smooth cycle, variable=\t]
      ({(1.18+0.09*abs(sin(6*\t)))*cos(\t)}, {0.1+(0.86+0.08*abs(sin(6*\t)))*sin(\t)});
  % hemisphere split and a few gyri squiggles
  \draw[#2, line width=0.8pt, line cap=round]
    (0,0.98) .. controls (0.12,0.8) and (-0.12,0.68) .. (0.02,0.52);
  \draw[#2, line width=0.7pt, line cap=round]
    (-0.85,0.45) .. controls (-0.75,0.7) and (-0.55,0.72) .. (-0.45,0.55)
                 .. controls (-0.38,0.42) and (-0.25,0.5) .. (-0.25,0.68);
  \draw[#2, line width=0.7pt, line cap=round]
    (0.3,0.72) .. controls (0.4,0.55) and (0.55,0.55) .. (0.6,0.7)
               .. controls (0.65,0.82) and (0.8,0.75) .. (0.88,0.55);
  \draw[#2, line width=0.7pt, line cap=round]
    (-1.0,0.0) .. controls (-0.9,0.15) and (-0.78,0.12) .. (-0.78,-0.05);
  \draw[#2, line width=0.7pt, line cap=round]
    (1.0,0.05) .. controls (0.92,0.2) and (0.8,0.18) .. (0.82,0.0);
}
\newcommand{\noggin@eyes}[3]{% #1 white radius, #2 pupil radius, #3 pupil shift (x,y)
  \foreach \ex in {-0.36,0.36}{
    \filldraw[fill=white, draw=ink, line width=0.8pt] (\ex,0.02) circle (#1);
    \fill[ink] ($(\ex,0.02)+(#3)$) circle (#2);
    \fill[white] ($(\ex,0.02)+(#3)+(0.035,0.04)$) circle (0.025);
  }
}
\newcommand{\noggin@cheeks}{%
  \fill[amRed!35, opacity=0.6] (-0.72,-0.25) ellipse (0.13 and 0.07);
  \fill[amRed!35, opacity=0.6] (0.72,-0.25) ellipse (0.13 and 0.07);
}
% --- faces --------------------------------------------------------------
\def\noggin@face@happy{%
  \noggin@body{brainpink}{brainline}\noggin@cheeks
  \noggin@eyes{0.2}{0.09}{0,-0.02}
  \draw[ink, line width=1pt, line cap=round] (-0.2,-0.32) .. controls (-0.1,-0.45) and (0.1,-0.45) .. (0.2,-0.32);
}
\def\noggin@face@panic{%
  \noggin@body{brainpink}{brainline}
  \noggin@eyes{0.25}{0.05}{0,0}
  \filldraw[fill=ink, draw=ink] (0,-0.42) ellipse (0.12 and 0.15);
  \fill[amBlue!60] (-1.05,0.55) .. controls (-1.15,0.4) and (-1.1,0.3) .. (-1.0,0.3)
                   .. controls (-0.9,0.3) and (-0.9,0.42) .. (-1.05,0.55);
  \fill[amBlue!60] (1.12,0.65) .. controls (1.02,0.5) and (1.07,0.4) .. (1.17,0.4)
                   .. controls (1.27,0.4) and (1.27,0.52) .. (1.12,0.65);
  \draw[amRed, line width=1.2pt, line cap=round] (-0.95,1.0) -- (-1.15,1.2);
  \draw[amRed, line width=1.2pt, line cap=round] (0,1.12) -- (0,1.38);
  \draw[amRed, line width=1.2pt, line cap=round] (0.95,1.0) -- (1.15,1.2);
}
\def\noggin@face@sleepy{%
  \noggin@body{brainpink}{brainline}\noggin@cheeks
  \foreach \ex in {-0.36,0.36}{
    \draw[ink, line width=1pt, line cap=round] (\ex-0.17,0.04) .. controls (\ex-0.08,-0.06) and (\ex+0.08,-0.06) .. (\ex+0.17,0.04);
  }
  \draw[ink, line width=0.9pt] (0.02,-0.38) circle (0.06);
  \node[font=\handfont\bfseries, text=amBlue, scale=0.9] at (1.15,1.0) {z};
  \node[font=\handfont\bfseries, text=amBlue, scale=1.3] at (1.45,1.3) {Z};
}
\def\noggin@face@confused{%
  \noggin@body{brainpink}{brainline}
  \noggin@eyes{0.2}{0.08}{0.06,0.07}
  \draw[ink, line width=1pt, line cap=round] (0.18,0.34) -- (0.52,0.42);
  \draw[ink, line width=1pt, line cap=round] (-0.52,0.28) -- (-0.2,0.28);
  \draw[ink, line width=1pt, line cap=round] (-0.2,-0.38) .. controls (-0.1,-0.3) and (-0.05,-0.46) .. (0.05,-0.38) .. controls (0.12,-0.31) and (0.17,-0.42) .. (0.22,-0.38);
  \node[font=\bfseries, text=amOrange!90!black, scale=1.8] at (1.2,1.05) {?};
}
\def\noggin@face@sick{%
  \noggin@body{sickgreen}{sickline}
  \foreach \ex in {-0.36,0.36}{
    \filldraw[fill=white, draw=ink, line width=0.8pt] (\ex,0.02) circle (0.2);
    \draw[ink, line width=0.8pt] plot[domain=0:900, samples=60, variable=\t]
      ({\ex+0.0002*\t*cos(\t)}, {0.02+0.0002*\t*sin(\t)});
  }
  \draw[ink, line width=1pt, line cap=round] (-0.24,-0.4) -- (-0.12,-0.33) -- (0,-0.4) -- (0.12,-0.33) -- (0.24,-0.4);
}
\def\noggin@face@hungry{%
  \noggin@body{brainpink}{brainline}\noggin@cheeks
  \noggin@eyes{0.22}{0.1}{0,0.04}
  \filldraw[fill=ink, draw=ink, line width=0.8pt] (-0.22,-0.3) .. controls (-0.15,-0.58) and (0.15,-0.58) .. (0.22,-0.3) -- cycle;
  \fill[amRed!60] (0.02,-0.47) ellipse (0.09 and 0.06);
  \fill[amBlue!45] (0.2,-0.48) .. controls (0.16,-0.6) and (0.18,-0.66) .. (0.23,-0.66)
                   .. controls (0.28,-0.66) and (0.3,-0.6) .. (0.2,-0.48);
}
\def\noggin@face@smug{%
  \noggin@body{brainpink}{brainline}\noggin@cheeks
  \noggin@eyes{0.2}{0.09}{0.05,-0.05}
  \foreach \ex in {-0.36,0.36}{
    \fill[brainpink] (\ex-0.21,0.02) rectangle (\ex+0.21,0.24);
    \draw[ink, line width=1pt] (\ex-0.2,0.04) -- (\ex+0.2,0.04);
  }
  \draw[ink, line width=1pt, line cap=round] (-0.18,-0.36) .. controls (0,-0.4) and (0.15,-0.38) .. (0.25,-0.26);
}
\tikzset{
  pics/noggin/.style={code={\csname noggin@face@#1\endcsname}},
}
\makeatother
\newcommand{\noggin}[2][1]{%
  \begin{tikzpicture}[scale=#1, baseline=(current bounding box.center)]
    \pic {noggin=#2};
  \end{tikzpicture}%
}

% TikZ + beamer overlays: \node[visible on=<2->] ... keeps the element in
% the layout (so nothing jumps) but hides it until slide 2.
\tikzset{
  invisible/.style={opacity=0, text opacity=0},
  visible on/.style={alt={#1{}{invisible}}},
  alt/.code args={<#1>#2#3}{\alt<#1>{\pgfkeysalso{#2}}{\pgfkeysalso{#3}}},
}

% Speech bubble for the mascot (handwritten font).
% \saybubble[<width>]{<text>}  -- use inside a tikzpicture via a node
\tikzset{
  say/.style={
    rectangle callout, draw=ink, line width=0.8pt, fill=white,
    rounded corners=4pt, inner sep=5pt, align=center,
    font=\handfont\normalsize, text=ink
  },
  think/.style={
    cloud callout, cloud puffs=12, draw=ink, line width=0.8pt, fill=white,
    inner sep=2pt, align=center, font=\handfont\normalsize, text=ink
  },
}

% =====================================================================
% Reusable boxes and frame building blocks
% =====================================================================
\tcbset{
  idiotbase/.style={
    enhanced, arc=3mm, boxrule=0pt, left=3mm, right=3mm, top=1.5mm,
    bottom=2mm, fonttitle=\bfseries\large, coltitle=white,
    attach boxed title to top left={xshift=4mm, yshift*=-2.5mm},
    boxed title style={arc=2mm, boxrule=0pt},
  }
}

% "Try it now" - a quick self-experiment students can do on the spot
\newtcolorbox{tryit}[1][]{idiotbase,
  colback=amOrange!12, colframe=amOrange, colbacktitle=amOrange!90!black,
  title={\faFlask\ \ Try it now}, borderline west={1.2mm}{0pt}{amOrange}, #1}

% "Poll" - for the webinar poll tool or a show of hands
\newtcolorbox{poll}[1][]{idiotbase,
  colback=amBlue!8, colframe=amBlue, colbacktitle=amBlue,
  title={\faPoll\ \ Poll}, borderline west={1.2mm}{0pt}{amBlue}, #1}

% "Chat" - students type into the chat
\newtcolorbox{chatbox}[1][]{idiotbase,
  colback=amBlue!8, colframe=amBlue, colbacktitle=amBlue,
  title={\faCommentDots\ \ Type in the chat}, borderline west={1.2mm}{0pt}{amBlue}, #1}

% "Science update" - where research moved on since the book (2016)
\newtcolorbox{scienceupdate}[1][]{idiotbase,
  colback=mist, colframe=ink, colbacktitle=ink, fonttitle=\bfseries\small,
  fontupper=\small, title={\faSyncAlt\ \ Science update}, #1}

% Plain rounded card
\newtcolorbox{card}[1][]{enhanced, arc=3mm, boxrule=0.8pt, colback=white,
  colframe=ink!25, left=3mm, right=3mm, top=2mm, bottom=2mm, #1}

% Myth-or-fact block (use inside a frame; the verdict appears on slide 2)
%   \mythorfact[<stamp colour>]{<statement>}{<verdict word>}{<explanation>}
\newcommand{\mythorfact}[4][amRed]{%
  \begin{center}
  \begin{tikzpicture}
    \node[draw=ink!30, line width=0.8pt, fill=white, rounded corners=4mm,
          inner xsep=8mm, inner ysep=6mm, text width=10.5cm, align=center,
          font=\Large\semiboldfont] (mf) {``#2''};
    \node[fill=ink, text=white, rounded corners=2mm, font=\bfseries\normalsize,
          inner sep=4pt, anchor=west] at ([xshift=6mm]mf.north west)
          {\faQuestionCircle\ \ MYTH or FACT?};
    \onslide<2->{
      \node[draw=#1, text=#1, line width=2.2pt, rounded corners=2mm, rotate=-12,
            font=\bfseries\LARGE, inner sep=5pt, fill=white, fill opacity=0.92,
            text opacity=1] at ([xshift=-6mm, yshift=1mm]mf.south east) {#3};
    }
  \end{tikzpicture}
  \end{center}
  \vspace{2mm}
  \onslide<2->{#4}%
}

% Big question: huge centred question with the confused mascot
%   \bigquestion{<question>}{<small line underneath>}
\newcommand{\bigquestion}[2]{%
  \vfill
  \begin{center}
  \begin{tikzpicture}
    \node[text width=10cm, align=center, font=\huge\bfseries, text=ink] (q) {#1};
    \node[below=4mm of q, text width=10cm, align=center, font=\large, text=softgray] {#2};
    \pic[scale=0.75] at ([xshift=-1.3cm, yshift=-0.3cm]q.west) {noggin=confused};
  \end{tikzpicture}
  \end{center}
  \vfill
}

% Title frame shared by all chapters. Set \title, \subtitle, \author,
% \date as usual, plus \chaptertagline{...}; then call \idiottitleframe{<mood>}.
\makeatletter
\newcommand{\chaptertagline}[1]{\def\idiot@tagline{#1}}
\def\idiot@tagline{}
\newcommand{\idiottitleframe}[1]{%
  \begin{frame}[plain]
    \begin{tikzpicture}[remember picture, overlay]
      \fill[paper] (current page.north west) rectangle (current page.south east);
      \fill[ink] (current page.north west) rectangle ([yshift=-1.2cm]current page.north east);
      \fill[amRed]    ([yshift=-1.2cm]current page.north west) rectangle ([yshift=-1.32cm, xshift=5.333cm]current page.north west);
      \fill[amBlue]   ([yshift=-1.2cm, xshift=5.333cm]current page.north west) rectangle ([yshift=-1.32cm, xshift=10.667cm]current page.north west);
      \fill[amOrange] ([yshift=-1.2cm, xshift=10.667cm]current page.north west) rectangle ([yshift=-1.32cm]current page.north east);
      \node[anchor=west, text=white, font=\large\bfseries]
        at ([xshift=8mm, yshift=-6mm]current page.north west) {\faBrain\ \ Book club for curious brains};
      \node[anchor=north west, text=ink, font=\fontsize{34}{38}\selectfont\bfseries]
        (t) at ([xshift=8mm, yshift=-2.1cm]current page.north west) {\inserttitle};
      \node[anchor=north west, text=amBlue, font=\LARGE\semiboldfont]
        (s) at ([yshift=-1mm]t.south west) {\insertsubtitle};
      \node[anchor=north west, text=ink, font=\large, text width=8.6cm]
        (g) at ([yshift=-3mm]s.south west) {\idiot@tagline};
      \node[anchor=south west, text=softgray, font=\normalsize, text width=9cm]
        at ([xshift=8mm, yshift=7mm]current page.south west) {\insertauthor\\ \insertdate};
      \fill[amOrange!18] ([xshift=-3.4cm, yshift=-0.6cm]current page.east) circle (2.6cm);
      \pic[scale=1.45] at ([xshift=-3.4cm, yshift=-0.4cm]current page.east) {noggin=#1};
    \end{tikzpicture}
  \end{frame}
}
\makeatother

% Section opener frame (creates the \section entry too).
%   \sectionopener{<number>}{<short name>}{<book's section title>}%
%                 {<what it is about>}{<fa icon>}{<colour>}{<mascot mood>}
\newcommand{\sectionopener}[7]{%
  \section{#2}%
  \begin{frame}[plain]
    \begin{tikzpicture}[remember picture, overlay]
      \fill[#6] (current page.north west) rectangle (current page.south east);
      \fill[white, opacity=0.08] ([xshift=-2cm,yshift=1cm]current page.north east) circle (5.5cm);
      \fill[white, opacity=0.06] ([xshift=1cm,yshift=-1cm]current page.south west) circle (3.5cm);
      \node[anchor=north west, text=white, font=\bfseries\Large, opacity=0.85]
        at ([xshift=9mm, yshift=-9mm]current page.north west) {Part #1 of 4};
      \node[anchor=west, text=white, font=\fontsize{30}{34}\selectfont\bfseries, text width=10cm]
        (st) at ([xshift=9mm, yshift=4mm]current page.west) {#3};
      \node[anchor=north west, text=white, font=\Large, text width=10cm, opacity=0.9]
        at ([yshift=-3mm]st.south west) {#4};
      \node[text=white, opacity=0.18, font=\fontsize{110}{110}\selectfont]
        at ([xshift=-3.2cm, yshift=0.6cm]current page.east) {#5};
      \pic[scale=0.9] at ([xshift=-3.0cm, yshift=-2.6cm]current page.east) {noggin=#7};
    \end{tikzpicture}
  \end{frame}
}
